% =====================================================================
% synopsis-genomic-scripting.tex
%
% A self-contained manuscript. It derives the theory it needs from
% first principles and cites only external literature. A reader is not
% expected to have read anything else by the authors.
% =====================================================================

\documentclass[11pt,a4paper]{article}

\usepackage[T1]{fontenc}
\usepackage[utf8]{inputenc}
\usepackage{lmodern}
\usepackage[margin=2.6cm]{geometry}
\usepackage{amsmath,amssymb,amsthm}
\usepackage{mathtools}
\usepackage{booktabs}
\usepackage{longtable}
\usepackage{array}
\usepackage{listings}
\usepackage{xcolor}
\usepackage{microtype}
\usepackage{enumitem}
\usepackage[hidelinks]{hyperref}
\usepackage[capitalise,nameinlink]{cleveref}
\usepackage{caption}

% ---------------------------------------------------------------------
% Theorem environments
% ---------------------------------------------------------------------
\theoremstyle{plain}
\newtheorem{theorem}{Theorem}[section]
\newtheorem{proposition}[theorem]{Proposition}
\newtheorem{lemma}[theorem]{Lemma}
\newtheorem{corollary}[theorem]{Corollary}

\theoremstyle{definition}
\newtheorem{definition}[theorem]{Definition}
\newtheorem{assumption}[theorem]{Assumption}
\newtheorem{construction}[theorem]{Construction}
\newtheorem{example}[theorem]{Example}
\newtheorem{rule}[theorem]{Rule}

\theoremstyle{remark}
\newtheorem{remark}[theorem]{Remark}
\newtheorem{convention}[theorem]{Convention}

\crefname{rule}{Rule}{Rules}
\Crefname{rule}{Rule}{Rules}
\crefname{assumption}{Assumption}{Assumptions}
\Crefname{assumption}{Assumption}{Assumptions}
\crefname{construction}{Construction}{Constructions}
\Crefname{construction}{Construction}{Constructions}

% ---------------------------------------------------------------------
% Notation
% ---------------------------------------------------------------------
\newcommand{\flo}{\beta}                       % contact floor
\newcommand{\sep}{\sigma}                      % separation cost
\newcommand{\med}{\mathfrak{m}}                % the medium
\newcommand{\Res}{\varrho}                     % residue
\newcommand{\Gph}{\mathcal{G}}
\newcommand{\Ver}{V}
\newcommand{\Edg}{E}
\newcommand{\weight}{w}
\newcommand{\Om}{\Omega}
\newcommand{\dem}{D}
\newcommand{\Rp}{\mathbb{R}_{>0}}
\newcommand{\Rnn}{\mathbb{R}_{\ge 0}}
\newcommand{\corr}{\varphi}
\newcommand{\resp}{\pi}
\newcommand{\typ}[1]{\ensuremath{\mathsf{#1}}}
\newcommand{\kw}[1]{\texttt{\textbf{#1}}}
\DeclareMathOperator{\cut}{cut}
\DeclareMathOperator*{\argmin}{arg\,min}

% ---------------------------------------------------------------------
% Listings: the synopsis language
% ---------------------------------------------------------------------
\definecolor{synkw}{HTML}{1F6FEB}
\definecolor{syntype}{HTML}{8250DF}
\definecolor{syncom}{HTML}{6E7781}
\definecolor{synstr}{HTML}{0A7F3F}
\definecolor{synerr}{HTML}{D1242F}
\definecolor{synbg}{HTML}{FAFAFA}

\lstdefinelanguage{synopsis}{
  morekeywords={under,let,seq,net,unit,response,coord,profile,ranked,
    peaks,verdict,real,depth,corr,claim,record,report,to,by,against,
    align,central,compare,detect,project,relax,until,quiescent,sweep,
    step,for,in,where,require,drop,bind,with,anchors,residue,medium,
    perturb,join,witness,frame,decline,correspond,diverge,assert,
    method,threshold,as,from,of,and,or,not,if,then,else,open,close},
  morekeywords=[2]{xcorr,spectral,shader,contact,cardinal,dna,protein,
    smith_waterman,jaccard,zscore,cosine,hierarchical},
  sensitive=true,
  morecomment=[l]{\#},
  morestring=[b]",
}

\lstset{
  language=synopsis,
  basicstyle=\ttfamily\small,
  keywordstyle=\color{synkw}\bfseries,
  keywordstyle=[2]\color{syntype},
  commentstyle=\color{syncom}\itshape,
  stringstyle=\color{synstr},
  backgroundcolor=\color{synbg},
  frame=leftline,
  framerule=1.2pt,
  rulecolor=\color{synkw!35},
  xleftmargin=1.6em,
  framexleftmargin=1.6em,
  showstringspaces=false,
  breaklines=true,
  breakatwhitespace=true,
  columns=fullflexible,
  numbers=left,
  numberstyle=\tiny\color{syncom},
  numbersep=8pt,
  escapeinside={(*@}{@*)},
  captionpos=b,
}

\lstdefinestyle{nonum}{numbers=none}
\lstdefinestyle{grammar}{
  language={},
  basicstyle=\ttfamily\footnotesize,
  keywordstyle={},
  numbers=none,
  frame=none,
  xleftmargin=1.2em,
  backgroundcolor=\color{white},
}

\hypersetup{
  colorlinks=false,
  pdftitle={Synopsis: a scripting language for sufficient alignment in genomics},
}

% =====================================================================

\title{\bfseries Synopsis: A Scripting Language for \\[0.2em]
       Sufficient Alignment in Genomics}

\author{}
\date{}

\begin{document}

\maketitle

\begin{abstract}
\noindent
Genomic comparison is conducted almost entirely in the vocabulary of
\emph{exact} correspondence: percent identity, edit distance, alignment
score. This vocabulary is not merely imprecise for the question
practitioners actually ask --- \emph{do these two things do the same
thing?} --- it is provably unable to answer it. We prove two separation
results. There exist pairs of units whose content is separated at the
theoretical minimum, so that no content-based criterion can distinguish
them, yet whose effect on the surrounding network diverges without
bound; and there exist pairs whose content is arbitrarily expensive to
reconcile yet whose effect is identical. Content comparison collapses
the first class to \emph{same} and the second to \emph{different}, and
is wrong in both directions. The two classes are orthogonal axes of a
plane that a scalar criterion projects onto a line.

We construct the theory needed to occupy that plane. Working in finite
weighted graphs with strictly positive weights, we establish a
\emph{contact floor}: a strictly positive lower bound $\flo$ on the cost
of any distinction. We show that identity is borne by a minimum cut ---
a region of edges --- and never by a vertex or a label, so that it
survives arbitrary relabelling exactly and is not carried by any name a
database can store. We show that the value a unit registers is a
function of the pair (unit, receiver) and not of the unit, and that the
resulting disagreement between receivers is not measurement error. We
then give a bidirectional comparison, \emph{cross-demand}, that
certifies correspondence without reading meaning, and prove that its
relaxation terminates in at most $\lceil \dem_0/\eta\rceil$ steps or
declines --- there is no third outcome and no divergence.

From this theory we design and specify \textbf{synopsis} (extension
\texttt{.syp}), a total scripting language for genomic experiments. Its
name is taken from the synoptic gospels, the books that are read in
parallel columns because they \emph{align sufficiently}, not because they
agree word for word; the discipline of reading them is exactly the
discipline we formalise. Synopsis is designed around refusals. Sequences
are opaque and cannot be indexed. Every value carries a residue bounded
below by the floor, and no constructor produces zero residue, so
\texttt{identical} is not expressible. The verdict constructor is
\emph{4-ary}: omitting the response columns is an arity error, not a
weaker result. Comparison across receiver frames is a scope error.
Thresholds have no defaults. Iteration is bounded by construction and
relaxation typechecks only when its step floor is bounded below, in which
case the compiler emits the termination bound as a proof obligation it has
discharged. We prove strong normalisation, floor preservation, scope
safety, and four \emph{refusal theorems} stating that certain scientifically
invalid claims have no well-typed representation.

We give the full grammar, a typing judgement, a small-step operational
semantics, and worked programs reproducing three families of genomic
experiment --- matched-filter motif detection, spectral homology search
over a database, and four-column functional adjudication --- in a single
language whose output is a report bearing certificates rather than a
number. We report honestly that one assumption the four-column verdict
requires, the independence of the verdict from the choice of admissible
response, \emph{fails} on solved biochemical networks: only $28.6\%$ of
comparisons fall within threshold, with a maximum discrepancy of $0.307$
against $\theta=0.01$. Synopsis makes this failure syntactically
visible: the response operator requires its map to be named, so no
program can conceal which arbitrary choice it made. We regard designing
a language that cannot hide its open problems as the point.
\end{abstract}

\tableofcontents

\clearpage

% =====================================================================
\section{Introduction}
\label{sec:intro}
% =====================================================================

\subsection{Two failures of the same kind}

Consider two questions a working geneticist asks daily.

\begin{enumerate}[label=(\roman*),leftmargin=2.2em]
\item \textbf{A synonymous variant.} A single base changes; the codon
      changes; the amino acid does not. Every content-based comparison
      reports the protein product as identical. Yet synonymous variants
      demonstrably alter translation kinetics, co-translational folding,
      mRNA stability and substrate specificity
      \cite{plotkin2011synonymous,sauna2011silent,kimchi2007silent}. The
      comparison that reports \emph{same} is not approximately right. It
      is reporting on a quantity that was never the one at issue.

\item \textbf{Two unrelated enzymes catalysing one reaction.} Non-homologous
      isofunctional enzymes are common enough to have a name and a
      systematic census \cite{omelchenko2010nonhomologous,galperin2010divergence}.
      Their sequences share nothing; every alignment reports maximal
      divergence. They do the same job. The comparison that reports
      \emph{different} is, again, answering a different question than
      the one asked.
\end{enumerate}

These are usually filed as separate difficulties --- one a false
positive, one a false negative --- to be patched by better substitution
matrices, deeper profile models, or structural prediction
\cite{henikoff1992amino,eddy2011accelerated,jumper2021highly}. We will
show they are one difficulty, and that it is structural rather than
statistical.

The claim we prove is sharp. Let \emph{content divergence} be any
quantity computed from the two units in isolation, and let
\emph{response divergence} be a quantity computed from what each unit
does to the network it sits in. \Cref{thm:falsefriend} constructs a
family in which content divergence is pinned at exactly zero --- not
small, exactly zero, at the theoretical minimum --- while response
divergence grows without bound. \Cref{cor:converse} constructs a family
in which response divergence is exactly zero while content divergence
grows. The two families lie on orthogonal axes of a two-dimensional
plane. Any criterion returning a scalar computed from content alone
projects that plane onto one axis, and necessarily maps one entire
family to the origin. No refinement of the scalar repairs this; the
defect is one of dimension, not of accuracy.

\subsection{Sufficient alignment}
\label{sec:sufficient}

The response to a comparison that cannot be exact is not to abandon
comparison. It is to change what one asks of it.

There is an old and instructive precedent. The first three canonical
gospels are called \emph{synoptic} --- seen together --- because they
can be laid out in parallel columns and read across
\cite{griesbach1776synopsis,throckmorton1949gospel}. They are not
identical. They differ in wording, in order, in detail, and in what each
chooses to omit. A concordance demanding word-for-word agreement finds
them in constant contradiction and produces nothing. The synoptic
reading demands something weaker and far more useful: that the columns
correspond \emph{sufficiently} for a passage in one to be located
against a passage in another, with the residue --- what does not
correspond --- retained and displayed rather than discarded. The
discipline is not to eliminate the difference. It is to bound it, to
know where it lies, and to be able to say when it is small enough not to
matter for the question at hand.

That is precisely the discipline genomic comparison needs and does not
have. We take the name of our language from it. A \emph{synopsis} is a
parallel-column reading in which correspondence is certified, residue is
carried, and the verdict is about sufficiency rather than identity.

The choice of name commits us to three design consequences, which we
discharge in \Cref{sec:design} and prove in \Cref{sec:meta}.

\begin{enumerate}[label=(\alph*),leftmargin=2.2em]
\item \textbf{Residue is never discarded.} A comparison that returns a
      single number has thrown the residue away. In synopsis every value
      carries its residue in its type, and there is no constructor
      producing zero residue (\Cref{thm:nozero}).
\item \textbf{Exactness is not expressible.} If \texttt{identical} can be
      written, it will be written, and the two failures of
      \Cref{sec:intro} return. In synopsis it has no constructor
      (\Cref{thm:noexact}).
\item \textbf{The columns are structural.} A synopsis with one column is
      not a synopsis. The verdict constructor of synopsis takes four
      columns and omitting a pair is an arity error, not a permitted
      simplification (\Cref{thm:arity}).
\end{enumerate}

\subsection{Where the invariant is not}
\label{sec:notinseq}

Before building, we should say plainly what we take to be the source of
the trouble, because it determines everything about the language's
treatment of sequence.

The tacit premise of content comparison is that whatever makes two units
functionally the same is \emph{written in them}, so that a sufficiently
clever reading of the two strings recovers it. We claim this premise is
false, and we can make the claim precise rather than rhetorical.

\begin{remark}[The invariant is not borne by the string]
\label{rem:bearing}
\Cref{thm:invariant-cut} shows that the quantity preserved under
relabelling of a contact structure is the weight of a minimum cut, and
\Cref{thm:region} shows that this cut is in general a region of edges
containing several vertices, not a single vertex. \Cref{cor:labels}
then shows that vertex labels may be permuted arbitrarily without
changing any separation cost. A sequence is a labelling. It is
therefore, by \Cref{cor:labels}, precisely the part of the structure
that the invariant is indifferent to. Reading harder does not help
because there is nothing further in the reading.
\end{remark}

There is an elementary observation in the same direction that requires no
theory at all. Every double-stranded nucleic acid carries two strands
\cite{watson1953molecular}. Read the complement in its own $5'\to3'$
direction and one obtains a different string, with a different codon
decomposition, and different open reading frames. Both strands are
physically present; both are read by the cell in different contexts;
neither is privileged by the chemistry. If the functional invariant were
a property of the string, the molecule would carry two of them and they
would disagree. What the molecule actually determines is the
\emph{pairing} --- the contact structure --- and both strings are
labellings of it. The invariant, if there is one, lives in the contacts.

Accordingly, in synopsis a sequence is an \textbf{opaque} type. It has no
indexing operator, no substring operator, no base accessor, and no
equality. It enters computation only by projection into a structure that
is a legitimate bearer of the invariant. This is the single most
restrictive decision in the language and \Cref{sec:opaque} defends it in
detail.

\subsection{Why a language, and why a total one}
\label{sec:whylanguage}

One could implement all of this as a library. We give three reasons not
to.

\textbf{First, a library cannot refuse.} The results of
\Cref{sec:separation} say that certain comparisons are invalid ---
comparing content alone, comparing across incompatible receiver frames,
reporting exact identity. In a library these are discouraged by
documentation. In a language they can be made unrepresentable. That is
the difference between a convention and a theorem, and
\Cref{sec:meta} proves four theorems of exactly this shape. A refusal
that the typechecker enforces is worth more than a paragraph in a manual
that a hurried user does not read.

\textbf{Second, defaults are where science goes wrong.} A library
function with a default threshold is an experiment with an
undeclared parameter. Multiplied across a pipeline this is the
mechanism by which analytic flexibility silently manufactures results
\cite{simmons2011false,gelman2013garden}. Synopsis has \emph{no default
thresholds anywhere}: every numeric parameter that can change a verdict
must be written in the program text, and the report prints it. This
costs the user keystrokes and buys the reader the ability to see what
was chosen.

\textbf{Third, an experiment should terminate.} A laboratory protocol
that might not finish is not a protocol. Synopsis is \emph{total}: every
well-typed program halts, and we prove it (\Cref{thm:total}). This is
purchased by giving up general recursion and unbounded loops
\cite{turner1995total}, which turns out to cost nothing --- none of the
genomic procedures we survey needs them --- and by requiring that the
one genuinely iterative construct, relaxation, be given a step floor, at
which point the compiler emits the iteration bound as a discharged
obligation (\Cref{sec:relax}).

This last point drives the one design decision we should flag explicitly
because a reader may expect otherwise. Synopsis is an \textbf{analytic}
language, not a generative one. It evaluates and certifies claims about
given material. It does \emph{not} search a space for sequences meeting a
target specification. Generative search requires a solver whose
termination depends on the geometry of the search space and cannot be
established by typing; admitting it would forfeit \Cref{thm:total},
which is the strongest property the language has. We judge the trade
correct: a total language that certifies is more useful to an
experimentalist than a partial one that also proposes, and the
generative layer, if wanted, can be built above synopsis and made to
emit synopsis programs for certification.

\subsection{Contributions and organisation}

\begin{itemize}[leftmargin=1.4em]
\item \Cref{sec:contact} develops contact structures and proves the
      contact floor (\Cref{thm:floor}), the invariance of cut weight
      (\Cref{thm:invariant-cut}), and its conditional survival under
      expansion (\Cref{thm:invariance}), with an explicit counterexample
      (\Cref{rem:not-positive}) showing the unconditional statement to
      be \emph{false}.
\item \Cref{sec:identity} proves that identity is regional
      (\Cref{thm:region}), that labels are free (\Cref{cor:labels}),
      and that registered value is receiver-relative
      (\Cref{thm:receiver}) with disagreement that is not noise
      (\Cref{cor:notnoise}).
\item \Cref{sec:blind} builds cross-demand and relaxation and proves
      monotonicity (\Cref{thm:mono}), the termination dichotomy
      (\Cref{thm:dichotomy}) and the form of quiescence
      (\Cref{thm:form}).
\item \Cref{sec:separation} proves the two separation results
      (\Cref{thm:falsefriend}, \Cref{cor:converse}) and derives the
      four-column criterion (\Cref{thm:four}).
\item \Cref{sec:design} states the eleven design rules of synopsis and
      argues each from a result above.
\item \Cref{sec:lang} gives the concrete language: lexis, grammar, types
      and the typing judgement.
\item \Cref{sec:semantics} gives the operational semantics.
\item \Cref{sec:meta} proves the metatheory: totality, floor
      preservation, scope safety, and the four refusal theorems.
\item \Cref{sec:programs} gives worked \texttt{.syp} programs for three
      families of experiment.
\item \Cref{sec:report} specifies the report and its certificates.
\item \Cref{sec:open} reports the negative result on response
      independence and what the language does about it.
\end{itemize}

% =====================================================================
\section{Contact structures and the floor}
\label{sec:contact}
% =====================================================================

We build the theory from nothing. A reader who wants only the language
may read \Cref{sec:design} onwards and refer back, but the design rules
are consequences of these results and are arbitrary without them.

\subsection{Finite weighted graphs}

\begin{definition}[Finite weighted graph]
\label{def:fwg}
A \emph{finite weighted graph} is a triple $\Gph=(\Ver,\Edg,\weight)$
with $\Ver$ a finite non-empty set, $\Edg\subseteq\binom{\Ver}{2}$, and
$\weight:\Edg\to\Rp$. Weights are \emph{strictly} positive: an edge of
weight zero is not a light edge, it is the absence of an edge, and we
insist on the distinction because every result below turns on it.
\end{definition}

\begin{definition}[Cut]
\label{def:cut}
For $S\subseteq\Ver$ the \emph{cut} induced by $S$ is
$\partial S=\{e\in\Edg: |e\cap S|=1\}$ with weight
$\cut(S)=\sum_{e\in\partial S}\weight(e)$. For $u\ne v$ a set $S$ is
\emph{$(u,v)$-separating} if $u\in S$ and $v\notin S$, and
\[
  \sep(u\mid v)\;=\;\min\{\cut(S)\;:\;S \text{ is }(u,v)\text{-separating}\}.
\]
The minimum is over a finite non-empty collection, hence attained.
\end{definition}

We call $\sep(u\mid v)$ the \emph{separation cost}: the cheapest way, in
the currency the weights supply, to place $u$ and $v$ on opposite sides
of a boundary. By the max-flow min-cut theorem it equals the maximum
flow from $u$ to $v$ \cite{ford1956maximal}, and is computable in
polynomial time \cite{stoer1997simple,karger1996new}; all-pairs values
admit a Gomory--Hu tree \cite{gomory1961multi}. We will not need
efficiency for the theory, but the language's cost model does
(\Cref{sec:costmodel}).

\begin{definition}[Contact structure]
\label{def:contact}
A \emph{contact structure} is a finite weighted graph with a
distinguished vertex $\med\in\Ver$, the \emph{medium}, adjacent to every
other vertex. Vertices other than $\med$ are \emph{items}, written
$I=\Ver\setminus\{\med\}$. We write
$\sep(v)=\sep(v\mid\med)$ for the \emph{separation cost of an item},
$\Om=\sum_{e\in\Edg}\weight(e)$ for the total weight, and
\[
  \flo(\Gph)\;=\;\min_{e\in\Edg}\weight(e)
\]
for the \emph{contact floor}.
\end{definition}

The medium formalises the background against which anything is
individuated. Every distinction is a distinction \emph{from} something;
$\med$ is that something, and its universal adjacency says the
background is always available and never absent. This is the graph-level
statement of a very old idea --- that a thing is known by what it is not
--- and it is what makes $\sep(v)$ the cost of individuating $v$ at all,
rather than merely of separating it from a chosen rival.

\subsection{The floor}

\begin{theorem}[Contact floor]
\label{thm:floor}
Let $\Gph$ be a contact structure and $v$ an item. Then
\[
  \sep(v)\;\ge\;\flo(\Gph)\;>\;0.
\]
\end{theorem}

\begin{proof}
$\Edg$ is finite and non-empty (the medium is adjacent to every item and
$I\neq\emptyset$), and $\weight$ takes values in $\Rp$, so
$\flo(\Gph)=\min_{e}\weight(e)$ is the minimum of a finite set of
strictly positive reals and is therefore strictly positive.

Let $S$ be any $(v,\med)$-separating set. Since $v\in S$ and
$\med\notin S$ and $\{v,\med\}\in\Edg$ by adjacency of the medium, we
have $|\{v,\med\}\cap S|=1$, so $\{v,\med\}\in\partial S$. Hence
$\partial S\neq\emptyset$ and
$\cut(S)\ge\weight(\{v,\med\})\ge\flo(\Gph)$. Taking the minimum over
separating $S$ gives $\sep(v)\ge\flo(\Gph)$.
\end{proof}

The proof is three lines and the content is substantial: \emph{no
distinction is free}. There is a strictly positive price on telling any
two things apart, fixed by the cheapest contact in the structure. Two
items may be as similar as the structure permits --- they may sit
exactly at the floor --- but they cannot be at zero, because zero is not
in the range of $\weight$.

This is the formal reason a comparison should never return
\texttt{identical}. Identity in the sense of \emph{indistinguishable at
cost zero} does not occur in a contact structure. Reporting it is
reporting a value outside the range of the quantity being measured.

\begin{definition}[Local floor]
\label{def:localfloor}
For an item $v$, $\flo_k(v)=\min\{\weight(e): v\in e\}$, the minimum
weight incident to $v$.
\end{definition}

Clearly $\flo(\Gph)\le\flo_k(v)$ for every $v$, and
$\flo(\Gph)=\min_v \flo_k(v)$.

\subsection{Invariance under relabelling}

\begin{definition}[Weighted isomorphism]
\label{def:iso}
A \emph{weighted isomorphism} of contact structures
$\Gph\to\Gph'$ is a bijection $h:\Ver\to\Ver'$ with $h(\med)=\med'$,
$\{h(u),h(v)\}\in\Edg' \iff \{u,v\}\in\Edg$, and
$\weight'(\{h(u),h(v)\})=\weight(\{u,v\})$.
\end{definition}

\begin{theorem}[Cut weight is the invariant]
\label{thm:invariant-cut}
If $h:\Gph\to\Gph'$ is a weighted isomorphism then for all items $u,v$,
\[
  \sep'(h(u)\mid h(v)) \;=\; \sep(u\mid v),
  \qquad
  \sep'(h(v)) \;=\; \sep(v),
  \qquad
  \flo(\Gph')=\flo(\Gph).
\]
\end{theorem}

\begin{proof}
$h$ induces a bijection $S\mapsto h(S)$ between subsets of $\Ver$ and of
$\Ver'$, carrying $(u,v)$-separating sets to $(h(u),h(v))$-separating
sets bijectively. It carries $\partial S$ to $\partial h(S)$ edge for
edge and preserves each weight, so $\cut'(h(S))=\cut(S)$. Minimising
over a bijectively-corresponding collection of equal values gives the
first identity; the second is the case $v=\med$, using $h(\med)=\med'$;
the third holds because $h$ is a weight-preserving bijection on $\Edg$.
\end{proof}

\begin{corollary}[Labels are free]
\label{cor:labels}
Let $\Gph$ be a contact structure and $\rho$ any permutation of the
items. Let $\Gph^\rho$ be $\Gph$ with items renamed by $\rho$. Then
every separation cost, every minimum cut weight and the floor are
unchanged.
\end{corollary}

\begin{proof}
$\rho$ extended by $\med\mapsto\med$ is a weighted isomorphism
$\Gph\to\Gph^\rho$. Apply \Cref{thm:invariant-cut}.
\end{proof}

\Cref{cor:labels} is the formal content of \Cref{rem:bearing}. The names
carry \emph{no} part of the invariant --- not a small part, none. A
comparison procedure that reads names is reading the one component of
the structure guaranteed to be irrelevant.

\subsection{Expansion, and a false theorem made true}

Real structures are never observed complete. One always has a subgraph
and may later add more. We must ask whether the floor survives.

\begin{definition}[Expansion]
\label{def:expansion}
An \emph{expansion} of $\Gph$ is a contact structure $\Gph^+$ with
$\Ver\subseteq\Ver^+$, $\Edg\subseteq\Edg^+$, $\med^+=\med$ and
$\weight^+|_\Edg=\weight$. An expansion \emph{sequence} is a chain
$\Gph_0\subseteq\Gph_1\subseteq\cdots$.
\end{definition}

It is tempting to assert that separation costs stay bounded below along
any expansion sequence. They do not.

\begin{remark}[The unconditional claim is false]
\label{rem:not-positive}
Fix an item $v$ and let $\Gph_k$ add to $\Gph_{k-1}$ a fresh vertex
$x_k$ joined to $v$ by an edge of weight $1/k$ (and to $\med$ at weight
$1$). Then $\flo(\Gph_k)\le 1/k\to 0$, and the local floor
$\flo_k(v)\to 0$. With weights $1/k^2$ the sum $\sum_k 1/k^2$ converges,
so one may arrange that $\sep(v)$ itself is driven towards a limit
determined entirely by the new light edges: $v$ \emph{dissolves into the
medium}. In the limit no positive lower bound survives. Any statement
that the floor is unconditionally preserved under expansion is therefore
false, and we record the counterexample rather than the temptation.
\end{remark}

The correct statement needs a hypothesis, and naming it precisely is the
point.

\begin{definition}[Floor-stabilising expansion]
\label{def:stabilising}
An expansion sequence $(\Gph_k)$ is \emph{floor-stabilising} at $v$ if
there is $\flo_\ast>0$ with $\flo_k(v)\ge\flo_\ast$ for all $k$: the
newly added contacts at $v$ do not become arbitrarily cheap.
\end{definition}

\begin{theorem}[Conditional invariance under expansion]
\label{thm:invariance}
Let $(\Gph_k)$ be an expansion sequence that is floor-stabilising at $v$
with constant $\flo_\ast$. Then $\sep_k(v)\ge\flo_\ast>0$ for all $k$,
and the sequence $(\sep_k(v))$ is bounded below by a positive constant
independent of $k$.
\end{theorem}

\begin{proof}
Fix $k$ and let $S$ be $(v,\med)$-separating in $\Gph_k$. As in
\Cref{thm:floor}, $\{v,\med\}\in\partial S$, so
$\cut_k(S)\ge\weight_k(\{v,\med\})\ge \flo_k(v)\ge\flo_\ast$. Minimise
over $S$.
\end{proof}

\begin{remark}
\Cref{thm:invariance} is weaker than one would like and exactly as
strong as is true. The hypothesis is checkable --- it is a statement
about the incident weights actually observed --- and a language can
require the user to assert it and can record the assertion in the
report. Synopsis does: see \Cref{sec:frames} and the \kw{require} form.
\end{remark}

\subsection{Perturbation}

\begin{definition}[Perturbation]
\label{def:perturbation}
A \emph{perturbation} of $\Gph$ is a map $\delta:\Edg\to\Rnn$; it acts by
$\weight\mapsto\weight+\delta$. Perturbations are non-negative:
they raise the cost of distinctions and never create separation for
free.
\end{definition}

\begin{proposition}[Skeleton invariance]
\label{prop:skeleton}
Under any perturbation, $\Ver$ and $\Edg$ are unchanged, every
separation cost is non-decreasing, and if $\delta$ vanishes on a
minimum-weight edge then $\flo$ is unchanged.
\end{proposition}

\begin{proof}
The vertex and edge sets do not appear in \Cref{def:perturbation}, so
they are fixed. For any $S$, $\cut_{\weight+\delta}(S)=\cut_\weight(S)+
\sum_{e\in\partial S}\delta(e)\ge \cut_\weight(S)$; minimising both
sides over the same finite collection preserves the inequality. If
$\delta(e_0)=0$ for some $e_0$ attaining the minimum weight, then
$\min_e(\weight+\delta)(e)\le \weight(e_0)=\flo$, and $\ge$ holds since
$\delta\ge0$.
\end{proof}

The distinction matters and is easy to lose: perturbation changes the
\emph{costs} without changing the \emph{contacts}. A unit which
perturbs its surroundings does not rewire them; it makes some
distinctions in them dearer. This is the sense in which a variant has an
effect without editing the network.

% =====================================================================
\section{Identity, region, and the receiver}
\label{sec:identity}
% =====================================================================

\subsection{Identity is a region, not a point}

\Cref{thm:invariant-cut} identifies the invariant as a cut weight. We
now show the minimising \emph{set} is generally not a singleton, so the
invariant is borne by a region.

\begin{theorem}[Regionality]
\label{thm:region}
There exist contact structures in which, for an item $v$, every
minimising $(v,\med)$-separating set has cardinality at least $3$.
\end{theorem}

\begin{proof}
Fix $r\ge3$ and $\flo>0$ with $\flo$ small. Take items
$A=\{a_1,\dots,a_r\}$ and $B=\{b_1,\dots,b_r\}$. Join every pair within
$A$ and every pair within $B$ by an edge of weight $W>0$; join $a_1$ to
$b_1$ by a single edge of weight $\flo$; join every item to $\med$ by an
edge of weight $\flo$. Let $v=a_1$.

Consider $S=\{a_1\}$. Its cut severs the $r-1$ edges from $a_1$ to the
rest of $A$, the bridge to $b_1$, and the medium edge, so
\[
  \cut(\{a_1\}) \;=\; (r-1)W + \flo + \flo \;=\; (r-1)W+2\flo .
\]
Now consider $S=A$. Its cut severs the bridge and the $r$ medium edges:
\[
  \cut(A) \;=\; \flo + r\flo \;=\; (r+1)\flo ,
\]
which is independent of $W$. Whenever $(r-1)W+2\flo > (r+1)\flo$, i.e.
$W>\flo\cdot\frac{r-1}{r-1}=\flo$, we get $\cut(A)<\cut(\{a_1\})$.
More generally, a separating set $S$ with $a_1\in S$ and
$|S\cap A|=j<r$ severs at least $j(r-j)$ intra-$A$ edges of weight $W$
plus $j$ medium edges, giving $\cut(S)\ge j(r-j)W+j\flo$, which for
$1\le j\le r-1$ is at least $(r-1)W+\flo$ and so exceeds $(r+1)\flo$
once $W>2\flo$. Taking $S\supseteq A$ and adding any $b_i$ only adds
intra-$B$ edges of weight $W$ and further medium edges, strictly
increasing the cut unless $S\supseteq A\cup B$, which is not separating
from $\med$'s side only if $\med\notin S$ --- and then
$\cut(A\cup B)=2r\flo>(r+1)\flo$ for $r>1$.

Hence for $W>2\flo$ and $r\ge3$ the unique minimiser is $S^\ast=A$, of
cardinality $r\ge3$.
\end{proof}

\begin{remark}[A witness we got wrong first]
\label{rem:badwitness}
Our first attempt at \Cref{thm:falsefriend} used the informal principle
``if two items are joined by an edge of weight $\flo$ then they are
separated at the floor''. This is false, and \Cref{thm:region} is why:
the minimum cut separating $a_1$ from $b_1$ in the construction above is
\emph{not} the bridge, because detaching the whole cluster $A$ may be
cheaper or dearer depending on $W$, and for moderate $W$ neither is the
bridge alone. Computing the cut exactly rather than reasoning from the
light edge gave $1.4$ where the informal argument predicted $0.2$. We
record the episode because the informal principle is intuitive, wrong,
and would have propagated silently into the language's semantics.
\end{remark}

\begin{corollary}
\label{cor:notpoint}
The invariant of \Cref{thm:invariant-cut} is not in general attributable
to any single vertex, and \emph{a fortiori} not to any single label.
\end{corollary}

\subsection{Receiver-relativity}

The value an item registers is not a property of the item.

\begin{definition}[Receiver]
\label{def:receiver}
A \emph{receiver} is a contact structure $N$ together with a designated
embedding site. A unit $u$ \emph{registers} in $N$ the value
$\sep_N(u)$, computed after $u$ is embedded at the site.
\end{definition}

\begin{theorem}[Receiver-relativity]
\label{thm:receiver}
There exist a unit $u$ and receivers $N,N'$ differing in a single edge
weight, with $\sep_N(u)\ne\sep_{N'}(u)$. Consequently $\sep(u)$ is a
function of the pair $(u,N)$ and not of $u$ alone.
\end{theorem}

\begin{proof}
Take $N$ with items $u,x,y$, edges $\{u,x\}$ and $\{x,y\}$ of weight $1$,
medium edges of weight $1$, and $\{u,y\}$ absent. Then a
$(u,\med)$-separating set must sever $\{u,\med\}$ and $\{u,x\}$ unless it
contains $x$, so $\sep_N(u)=\min\{1+1,\;\,1+1+1\}=2$ where the second
option is $S=\{u,x\}$ severing $\{u,\med\},\{x,\med\},\{x,y\}$. Now let
$N'$ differ only in $\weight(\{x,y\})=\tfrac14$. Then $S=\{u,x\}$ costs
$1+1+\tfrac14=2.25$ and $S=\{u\}$ still costs $2$, so $\sep_{N'}(u)=2$.
Instead set $\weight(\{u,x\})=\tfrac14$ in $N'$: then $S=\{u\}$ costs
$1+\tfrac14=1.25<2$, so $\sep_{N'}(u)=1.25\neq 2=\sep_N(u)$, and $N,N'$
differ in exactly one weight.
\end{proof}

\begin{corollary}[Disagreement is not noise]
\label{cor:notnoise}
Let $N,N'$ be receivers with $\sep_N(u)\ne\sep_{N'}(u)$. Then no
increase in the precision or number of measurements performed in either
receiver reduces the discrepancy, since each value is exactly determined
by its own structure. The two receivers are evaluating different
functions and each is correct.
\end{corollary}

\Cref{cor:notnoise} has an immediate practical reading and a language
consequence. The practical reading is that a variant reported with
different effect in two assays need not indicate that one assay is
wrong; the assays may be different receivers, both right. The language
consequence is that values from different receivers must not be
compared, and synopsis enforces this lexically: values are scoped to
their frame and cross-frame comparison is a \emph{scope error}, not a
runtime warning (\Cref{sec:frames}, \Cref{thm:scope}).

% =====================================================================
\section{Blind comparison: demand, relaxation, quiescence}
\label{sec:blind}
% =====================================================================

We now build a comparison that certifies correspondence without reading
what anything means. This is essential: any procedure that must read
meaning inherits the annotation errors already endemic in the databases
it reads \cite{schnoes2009annotation}.

\subsection{Alignment and distance to the floor}

\begin{definition}[Alignment]
\label{def:alignment}
For items $u,v$ of a contact structure $\Gph$,
\[
  a(u,v)\;=\;\frac{\sep(u\mid v)}{\Om}, \qquad
  \frac{\flo}{\Om}\ \text{ the \emph{floor ratio}}, \qquad
  \Delta(u,v)\;=\;a(u,v)-\frac{\flo}{\Om}.
\]
$\Delta$ is the \emph{distance to the floor}.
\end{definition}

\begin{proposition}
\label{prop:nonneg}
$\Delta(u,v)\ge0$ always, with equality only if $u$ and $v$ are
separated exactly at the floor.
\end{proposition}

\begin{proof}
Immediate from \Cref{thm:floor} applied to the pair $(u,v)$: any
separating set severs at least one edge, of weight at least $\flo$.
\end{proof}

Normalising by $\Om$ makes $a$ dimensionless and comparable across
structures of different total weight; subtracting the floor ratio makes
$\Delta$ vanish exactly when the pair is as close as the structure
permits. \emph{Sufficient} alignment will mean $\Delta$ small; it will
never mean $a=0$, which \Cref{thm:floor} forbids.

\subsection{Anchors and residue}

Not every item in a structure is equally load-bearing for a given item's
individuation. We split by ablation.

\begin{definition}[Ablation score, anchors, residue]
\label{def:anchor}
For items $x\ne u$, the \emph{ablation score} is
$\alpha(x;u)=\sep_\Gph(x)-\sep_{\Gph\setminus u}(x)$, the drop in $x$'s
separation cost when $u$ is deleted. Fix $r\ge0$. The \emph{anchors} of
$x$ are the $r$ items of largest $\alpha(x;\cdot)$, ties broken by a
fixed total order on names; the \emph{residue} $\Res(x)$ is the
remainder.
\end{definition}

The intuition is that anchors are what $x$ is held in place by, and
residue is everything else --- the part of the structure whose
correspondence still has to be established. The tie-breaking by a fixed
order is not cosmetic: without it the split is not a function, and a
language that intends to be reproducible cannot afford a
nondeterministic primitive. Synopsis therefore fixes the order and the
report prints $r$.

\subsection{Cross-demand}

\begin{definition}[Cross-demand]
\label{def:crossdemand}
Let $\Gph_A,\Gph_B$ be contact structures and
$\corr: I_B \rightharpoonup I_A$ a partial \emph{correspondence}. The
\emph{demand from $A$ on $B$} is
\begin{equation}
\label{eq:demand}
  \dem_{A\to B}
  \;=\;
  \sum_{i\in \Res_B}
  \max\!\left(0,\;
    a_B\!\left(i,\corr(i)\right)-\frac{\flo_B}{\Om_B}
  \right),
\end{equation}
the sum taken over residue items of $B$ on which $\corr$ is defined.
\end{definition}

Each term is a distance to the floor and is non-negative by
\Cref{prop:nonneg}; the clamp guards only floating-point noise. The
quantity is a \emph{demand} in the sense that it measures how much
reconciliation $B$ still requires of $A$: it is zero exactly when every
residue item of $B$ is at the floor from its correspondent.

Crucially, \eqref{eq:demand} reads no labels, no annotations, no
sequence. It is computed from cuts. That is what makes it blind.

\begin{definition}[Total demand]
$\dem = \dem_{A\to B} + \dem_{B\to A}$, using $\corr$ and its inverse.
\end{definition}

Bidirectionality is not decoration. A one-sided demand can be made
small by choosing a correspondence that maps everything difficult to
undefined; the reverse direction penalises exactly that.

\subsection{Relaxation}

\begin{definition}[Relaxation]
\label{def:relax}
Let $\dem_0\ge0$. A \emph{relaxation} is a sequence
$\dem_0>\dem_1>\cdots$ produced by an update $U$ with
$\dem_{k+1}=U(\dem_k)$, terminating when $\dem_k\le\theta$ for a stated
threshold $\theta\ge0$. The relaxation \emph{declines} at step $k$ if
$U(\dem_k)$ is undefined --- no effective update exists.
\end{definition}

\begin{assumption}[Effective update]
\label{ass:effective}
There is $\eta>0$ such that whenever $U(\dem)$ is defined,
$U(\dem)\le \dem-\eta$.
\end{assumption}

\Cref{ass:effective} is the whole of what makes relaxation a
terminating procedure, and it is a real assumption, not a technicality:
an update permitted to make arbitrarily small progress can run forever
while always improving. Synopsis therefore refuses to typecheck a
relaxation whose $\eta$ is not bounded below by a written constant
(\Cref{sec:relax}).

\begin{theorem}[Monotone decrease]
\label{thm:mono}
Under \Cref{ass:effective} the sequence $(\dem_k)$ is strictly
decreasing while defined, and $\dem_k\le \dem_0-k\eta$.
\end{theorem}

\begin{proof}
Induction: $\dem_0\le\dem_0-0$; if $\dem_k\le\dem_0-k\eta$ and
$U(\dem_k)$ is defined then
$\dem_{k+1}\le \dem_k-\eta\le \dem_0-(k+1)\eta$.
\end{proof}

\begin{theorem}[Termination dichotomy]
\label{thm:dichotomy}
Under \Cref{ass:effective}, exactly one of the following occurs within
$\lceil \dem_0/\eta\rceil$ steps:
\begin{enumerate}[label=(\roman*),leftmargin=2.2em]
\item \textbf{quiescence:} $\dem_k\le\theta$ for some
      $k\le\lceil \dem_0/\eta\rceil$; or
\item \textbf{decline:} $U(\dem_k)$ is undefined for some
      $k\le\lceil \dem_0/\eta\rceil$.
\end{enumerate}
There is no third outcome, and in particular the relaxation cannot
diverge or cycle.
\end{theorem}

\begin{proof}
Suppose neither (i) nor (ii) occurs for any
$k\le K=\lceil \dem_0/\eta\rceil$. Then $U$ is defined at every such
step and $\dem_k>\theta\ge0$ for all $k\le K$. By \Cref{thm:mono},
$\dem_K\le \dem_0-K\eta\le \dem_0-\dem_0=0$, contradicting
$\dem_K>\theta\ge0$. Exclusivity holds because (i) terminates the
sequence and (ii) is the failure of the update that would continue it.
\end{proof}

\Cref{thm:dichotomy} is the theorem that lets synopsis be total while
still containing a loop. The bound $\lceil \dem_0/\eta\rceil$ is
computable from values in scope at the point the loop is written, so the
compiler can emit it. \Cref{sec:relax} does exactly that.

\begin{theorem}[Form of quiescence]
\label{thm:form}
Quiescence at $\theta=0$ means every residue term of \eqref{eq:demand}
is zero, i.e.\ every corresponding residue pair is separated exactly at
the floor. It does \emph{not} mean any pair is indistinguishable: by
\Cref{thm:floor} each such pair still costs $\flo>0$ to separate.
\end{theorem}

\begin{proof}
A finite sum of non-negative terms is zero iff each term is zero; each
term is $\Delta(i,\corr(i))$, which vanishes iff the pair is at the
floor (\Cref{prop:nonneg}). The second statement is \Cref{thm:floor}.
\end{proof}

\Cref{thm:form} is the formal statement of \emph{sufficient alignment}.
Two structures may place no demand on one another and remain fully
distinguishable. This is the synoptic reading made precise: the columns
correspond; the residue is positive; both facts are reported.

% =====================================================================
\section{The separation results and the four-column criterion}
\label{sec:separation}
% =====================================================================

\subsection{Content-aligned, response-divergent}

\begin{theorem}[False friends]
\label{thm:falsefriend}
There is a family of pairs $(u_c,v_c)$, indexed by $c\ge0$, embedded in
contact structures, such that
\[
  \dem^{\mathrm{central}}(u_c,v_c)=0 \quad\text{for all } c,
  \qquad\text{while}\qquad
  \dem^{\mathrm{response}}(u_c,v_c)\longrightarrow \text{unbounded as } c\to\infty .
\]
That is: content is aligned exactly at the floor for every $c$, while
the effect on the surrounding network diverges.
\end{theorem}

\begin{proof}
\emph{Construction.} Let $P$ be a contact structure containing two
pendant items $u,v$, each joined to the medium by an edge of weight
$\flo$ and to nothing else. Then any $(u,v)$-separating set may take
$S=\{u\}$, severing exactly $\{u,\med\}$, so
$\sep(u\mid v)=\flo$ and $\Delta(u,v)=0$. Being pendant is essential;
by \Cref{rem:badwitness} a light edge between non-pendant items does
\emph{not} suffice.

Take $\Gph_A$ and $\Gph_B$ each equal to $P$, with $\corr$ the identity
on the pendants. All residue terms are $\Delta=0$, so
$\dem^{\mathrm{central}}=0$ for every $c$: content comparison reports
correspondence unconditionally.

\emph{Response.} Now let $u$ induce the perturbation $\delta_u\equiv0$
on a receiver $N$, and let $v$ induce $\delta_v$ raising a designated
edge $e_\ast$ of $N$ by $c$. By \Cref{prop:skeleton} the skeleton is
unchanged in both cases, so the response structures $N_u$ and $N_v$ are
comparable edge for edge under the identity correspondence. Choose
$e_\ast$ on the minimum cut of some residue item $i$; then
$\sep_{N_v}(i)=\sep_{N_u}(i)+c$ for $c$ small enough that the minimiser
does not change, and beyond that $\sep_{N_v}(i)$ continues to increase,
being the minimum of a family of affine functions of $c$ with at least
one coefficient positive. Hence $\Delta_{N_v}(i,\corr(i))$ grows without
bound in $c$ up to the normalisation $\Om$, which also grows linearly;
choosing $e_\ast$ so that its weight is a vanishing fraction of $\Om$
(by padding $N$ with fixed edges) makes the normalised gap grow, and
$\dem^{\mathrm{response}}\to\infty$ accordingly.
\end{proof}

\begin{remark}
The biological instance is (i) of \Cref{sec:intro}. Two units whose
content a comparison cannot separate --- a synonymous pair --- may act
on the cell in arbitrarily different ways
\cite{kimchi2007silent,sauna2011silent}. The theorem says this is not an
edge case to be handled but a structural possibility that content
comparison is blind to by construction.
\end{remark}

\subsection{Content-divergent, response-aligned}

\begin{corollary}[Converse]
\label{cor:converse}
There is a family indexed by $W$ with
$\dem^{\mathrm{central}}$ strictly increasing in $W$ and
$\dem^{\mathrm{response}}=0$ for all $W$.
\end{corollary}

\begin{proof}
Take the two-cluster structure of \Cref{thm:region} with $u=a_1$,
$v=b_1$. As $W$ grows, the minimum $(u,v)$-cut is $\min\{(r+1)\flo,\;
(r-1)W+2\flo,\dots\}$; for $W$ in the regime where a cluster cut is
selected the value is constant, but normalising by
$\Om=r(r-1)W+\flo+2r\flo$ makes $a(u,v)$ vary, and choosing the regime
$W<\flo$ where $S=\{a_1\}$ is minimal gives
$a(u,v)=\frac{(r-1)W+2\flo}{\Om}$, strictly increasing in $W$ over the
admissible interval. So $\dem^{\mathrm{central}}$ increases.

Let both $u$ and $v$ induce the \emph{same} perturbation $\delta$ on the
receiver. Then $N_u=N_v$ identically, every $\Delta$ term is computed on
one structure against itself under the identity correspondence, and
$\dem^{\mathrm{response}}=0$ exactly.
\end{proof}

\begin{remark}
This is (ii) of \Cref{sec:intro}: the non-homologous isofunctional
enzymes \cite{omelchenko2010nonhomologous}. Content comparison reports
maximal divergence; the response columns report correspondence.
\end{remark}

\subsection{The plane and the criterion}

\begin{corollary}[Orthogonality]
\label{cor:orth}
The families of \Cref{thm:falsefriend} and \Cref{cor:converse} lie on
the two coordinate axes of the plane
$(\dem^{\mathrm{central}},\dem^{\mathrm{response}})$: the first on
$\{0\}\times\Rnn$, the second on $\Rnn\times\{0\}$. Any criterion that
is a function of $\dem^{\mathrm{central}}$ alone is constant on the
first family and therefore cannot separate its members; any criterion
that is a function of $\dem^{\mathrm{response}}$ alone is constant on
the second.
\end{corollary}

This forces the shape of a correct criterion.

\begin{definition}[Four columns]
\label{def:fourq}
A \emph{four-column comparison} consists of two \emph{central} columns
$(C_A, C_B)$ --- the units under comparison, with their anchors and
residue --- and two \emph{response} columns $(R_A, R_B)$ --- the
receivers as perturbed by each unit --- together with correspondences
$\corr^{\mathrm{c}}$ and $\corr^{\mathrm{r}}$ and a threshold $\theta>0$.
Its verdict is
\[
\mathrm{verdict}=
\begin{cases}
\textsf{CORRESPOND} & \dem^{\mathrm{c}}\le\theta \text{ and } \dem^{\mathrm{r}}\le\theta,\\[2pt]
\textsf{DECLINE}    & \text{some relaxation declined (\Cref{thm:dichotomy}(ii))},\\[2pt]
\textsf{DIVERGE}    & \text{otherwise.}
\end{cases}
\]
\end{definition}

\begin{theorem}[Adequacy of four columns]
\label{thm:four}
The four-column verdict separates the families of
\Cref{thm:falsefriend} and \Cref{cor:converse}, and no criterion
depending on fewer than both demand coordinates does.
\end{theorem}

\begin{proof}
On the first family $\dem^{\mathrm{c}}=0\le\theta$ but
$\dem^{\mathrm{r}}$ exceeds $\theta$ for $c$ large, so the verdict is
\textsf{DIVERGE}, correctly. On the second $\dem^{\mathrm{r}}=0\le\theta$
but $\dem^{\mathrm{c}}$ exceeds $\theta$ for $W$ large, so again
\textsf{DIVERGE} --- and at small $W$, \textsf{CORRESPOND}, tracking the
genuine functional agreement. The impossibility half is
\Cref{cor:orth}.
\end{proof}

\begin{remark}[Why \textsf{DECLINE} is a verdict]
\label{rem:decline}
A three-valued verdict is not a hedge. \Cref{thm:dichotomy} gives
exactly two ways a relaxation can end, and one of them is the discovery
that no effective update exists --- the comparison, run correctly to its
end, does not decide. Collapsing that onto \textsf{DIVERGE} reports a
finding where there is none, and collapsing it onto \textsf{CORRESPOND}
is worse. On a laboratory bench ``this experiment does not decide'' is a
result one writes down. Synopsis makes \textsf{DECLINE} a first-class
inhabitant of \typ{verdict} for that reason.
\end{remark}

\subsection{The assumption that fails}

\Cref{def:fourq} requires, for each unit, a perturbation it induces ---
its \emph{admissible response}. Nothing so far fixes which one.

\begin{assumption}[Response independence]
\label{ass:respind}
The four-column verdict is independent of the choice of admissible
response, within threshold $\theta$.
\end{assumption}

\begin{proposition}
\label{prop:respind}
If \Cref{ass:respind} fails then the verdict depends on an arbitrary
choice made by the analyst, and two analysts running the same comparison
may correctly obtain different verdicts.
\end{proposition}

\begin{proof}
Immediate: the verdict is a function of $\dem^{\mathrm{r}}$, which is a
function of the chosen response.
\end{proof}

We report in \Cref{sec:open} that \Cref{ass:respind} \emph{fails
empirically} on solved biochemical networks. We do not repair it here.
We instead require the language to make the choice visible, which is the
most a language can honestly do about an open problem: see
\Cref{rule:namedresponse} and \Cref{sec:open}.

% =====================================================================
\section{Design: eleven rules}
\label{sec:design}
% =====================================================================

Every rule below is derived from a result above. We state the rule, the
result forcing it, and the consequence for a user.

\begin{rule}[Sequences are opaque]
\label{rule:opaque}
The type \typ{seq} supports no indexing, no slicing, no base access, no
concatenation and no equality. A \typ{seq} may only be consumed by a
projection into \typ{profile}, \typ{coord} or \typ{net}.
\end{rule}
\noindent\emph{Forced by} \Cref{cor:labels} and \Cref{rem:bearing}: a
sequence is a labelling, and labels carry none of the invariant.
\emph{Consequence:} a user cannot write a character-level comparison,
because there is no expression denoting a character. This is discussed
at length in \Cref{sec:opaque} because it is the rule users find hardest.

\begin{rule}[Every value carries residue]
\label{rule:residue}
Every value of a comparison type carries a pair $(\flo,\Res)$ with
$\Res\ge\flo>0$. Residue is not optional and not erasable; it may be
consumed by \kw{bind} or explicitly abandoned by \kw{drop}, which the
report records.
\end{rule}
\noindent\emph{Forced by} \Cref{thm:floor}. \emph{Consequence:} a program
that throws information away must say so in its own text.

\begin{rule}[No zero-residue constructor]
\label{rule:nozero}
There is no expression in synopsis whose type ascribes $\Res=0$.
\end{rule}
\noindent\emph{Forced by} \Cref{thm:floor}: zero is outside the range of
the quantity. \emph{Consequence:} \texttt{identical} and
\texttt{percent\_identity == 100} are not writable
(\Cref{thm:noexact}).

\begin{rule}[Verdicts are four-ary]
\label{rule:arity}
The verdict constructor \kw{align} takes exactly a central pair and a
response pair. Supplying only the central pair is an arity error at
parse time.
\end{rule}
\noindent\emph{Forced by} \Cref{cor:orth} and \Cref{thm:four}.
\emph{Consequence:} the commonest scientific error in this domain
--- concluding functional sameness from content sameness --- is a
\emph{syntax} error.

\begin{rule}[Frames are lexical]
\label{rule:frames}
Every value is created inside an \kw{under} block naming a receiver and
an encoding. Values from different blocks have different types and
cannot be compared.
\end{rule}
\noindent\emph{Forced by} \Cref{thm:receiver} and \Cref{cor:notnoise}.
\emph{Consequence:} comparing a nucleotide embedding to a protein
embedding, or a value from receiver $N$ to one from $N'$, fails to
typecheck (\Cref{thm:scope}).

\begin{rule}[No default thresholds]
\label{rule:nodefault}
Every numeric parameter that can change a verdict must appear in the
program text. There are no defaults.
\end{rule}
\noindent\emph{Forced by} nothing in the mathematics --- this one is
methodological \cite{simmons2011false,gelman2013garden}.
\emph{Consequence:} the report can print the full parameter set because
the parser has seen all of it.

\begin{rule}[Named responses]
\label{rule:namedresponse}
The response constructor is \kw{response}~$u$~\kw{by}~$\resp$, where
$\resp$ must be a named, in-scope perturbation map. There is no anonymous
or inferred response.
\end{rule}
\noindent\emph{Forced by} \Cref{prop:respind} and the negative result of
\Cref{sec:open}. \emph{Consequence:} every program records which
arbitrary choice it made, so the open problem is visible in the source
and in the report rather than buried in a library default.

\begin{rule}[Bounded iteration only]
\label{rule:bounded}
The only iteration forms are \kw{sweep} over an arithmetic range or a
literal list, \kw{for} over the finite item set of a structure, and
\kw{relax until quiescent}. There is no \kw{while}, no recursion, and no
user-defined fixpoint.
\end{rule}
\noindent\emph{Forced by} the desire for \Cref{thm:total}, made
affordable by \Cref{thm:dichotomy}. \emph{Consequence:} every synopsis
program halts, and the compiler can state the bound.

\begin{rule}[Relaxation must be floored]
\label{rule:floored}
\kw{relax}~\ldots~\kw{until quiescent} typechecks only when a step floor
$\eta>0$ is supplied as a literal or an expression provably positive.
The compiler then emits $\lceil \dem_0/\eta\rceil$ into the report.
\end{rule}
\noindent\emph{Forced by} \Cref{ass:effective} and \Cref{thm:dichotomy}.
\emph{Consequence:} the termination proof of each program is a
by-product of compiling it.

\begin{rule}[Three result types]
\label{rule:three}
Comparison methods produce distinct types according to their index set:
\typ{profile} (indexed by lag), \typ{ranked} (indexed by database
entry), \typ{verdict} (an element of a three-valued set over a
structure). These do not unify and are not coercible.
\end{rule}
\noindent\emph{Forced by} the fact that the index sets genuinely differ;
see \Cref{sec:threefamilies}. \emph{Consequence:} a user cannot
accidentally feed a lag-indexed correlation into a database-ranking
consumer.

\begin{rule}[Programs emit reports]
\label{rule:report}
The terminal form of a synopsis program is a report: a document carrying
claims, their verdicts, their residues, their parameters and their
certificates. A program does not ``return a number''.
\end{rule}
\noindent\emph{Forced by} \Cref{thm:form}: quiescence coexists with
positive residue, so a scalar cannot express the result.
\emph{Consequence:} outputs are auditable in the sense of
\cite{sandve2013ten,wilkinson2016fair,moreau2011provenance}.

\subsection{On \Cref{rule:opaque}, at length}
\label{sec:opaque}

Users object to opacity, so it deserves a defence beyond the citation of
\Cref{cor:labels}.

The objection is that sequences are data and data should be
manipulable. The answer is that in synopsis the sequence is not the
datum; the \emph{structure it labels} is. Opening a \typ{seq} to
indexing would immediately permit writing
\texttt{s[i] == t[i]} and, from there, percent identity --- which
\Cref{rule:nozero} and \Cref{thm:noexact} exist to forbid. A rule that
can be circumvented by three lines of user code is not a rule.

The practical question is then whether opacity costs anything. We claim
it does not, on the evidence of the three experiment families we
specify in \Cref{sec:programs}. Each begins by \emph{channelising} the
sequence --- mapping it to a numeric matrix under a declared encoding
--- and never touches a base again. Motif detection consumes channels
via correlation; homology search consumes them via a spectral transform;
functional adjudication consumes a contact structure derived from
pairing rather than from the string. None of them needs
\texttt{s[i]}.

What opacity does cost is the ability to write \emph{ad hoc} string
surgery inside a synopsis program. We regard this as a feature: such
surgery is data preparation, belongs upstream, and should be visible in
the provenance record rather than hidden inside an analysis script.

\subsection{Three families, three index sets}
\label{sec:threefamilies}

\Cref{rule:three} deserves justification because the temptation to unify
is strong and we ourselves initially yielded to it.

\begin{itemize}[leftmargin=1.4em]
\item \textbf{Matched filtering.} A query of length $L_q$ is slid along a
      target of length $L_t$ and a score computed at each offset. The
      result is indexed by \emph{lag} $k\in\{0,\dots,L_t-L_q\}$
      \cite{turin1960introduction,oppenheim2009discrete}. Its natural
      consumer is peak detection; its natural output is a set of
      \emph{positions}.
\item \textbf{Database search.} A query embedding is compared against
      $N$ stored embeddings. The result is indexed by \emph{database
      entry} $i\in\{0,\dots,N-1\}$
      \cite{altschul1990basic,steinegger2017mmseqs2}. Its natural
      consumer is top-$K$ selection; its natural output is a
      \emph{ranking}.
\item \textbf{Functional adjudication.} Two structures are compared. The
      result is not indexed at all: it is a single element of
      $\{\textsf{CORRESPOND},\textsf{DIVERGE},\textsf{DECLINE}\}$ with
      four demand values attached.
\end{itemize}

A single \texttt{compare} returning a single type would have to pick one
of these index sets and coerce the others. Coercing a lag index to a
database index is meaningless; coercing either to a scalar discards
exactly the information the method exists to produce. Synopsis therefore
keeps one \emph{keyword}, \kw{compare}, and three \emph{result types},
with the parser dispatching on the named method. The user writes the
same verb; the typechecker knows which noun came back.

% =====================================================================
\section{The synopsis language}
\label{sec:lang}
% =====================================================================

Files carry the extension \texttt{.syp}.

\subsection{Lexical structure}

Comments run from \texttt{\#} to end of line. Identifiers match
\texttt{[A-Za-z\_][A-Za-z0-9\_]*}. Numeric literals are decimal, with an
optional exponent; \emph{there are no unsuffixed bare numbers in
threshold position} --- see \Cref{sec:units}. String literals are
double-quoted and are used only for file paths and claim text. The
language is case-sensitive and whitespace-insensitive except that
newlines terminate statements.

\subsection{Types}

\begin{center}
\begin{tabular}{@{}llp{7.4cm}@{}}
\toprule
\textbf{Type} & \textbf{Carries} & \textbf{Meaning} \\
\midrule
\typ{seq}      & ---                 & An opaque sequence. No observers. \\
\typ{frame}    & receiver, encoding  & A lexical scope for values. \\
\typ{profile}$_\phi$ & $(\flo,\Res)$, lag index & Lag-indexed score vector in frame $\phi$. \\
\typ{coord}$_\phi$   & $(\flo,\Res)$, dimension  & A point in an embedding space. \\
\typ{ranked}$_\phi$  & $(\flo,\Res)$, entry index & A ranking over a database. \\
\typ{peaks}$_\phi$   & $(\flo,\Res)$, positions   & Detected positions with scores. \\
\typ{net}$_\phi$     & $(\flo,\Res)$, $\Om$       & A contact structure. \\
\typ{unit}$_\phi$    & $(\flo,\Res)$, anchors, residue & A central column. \\
\typ{response}$_\phi$ & $(\flo,\Res)$, named $\resp$ & A response column. \\
\typ{corr}$_{\phi,\psi}$ & partial map          & A correspondence between two structures. \\
\typ{verdict}  & four demands, $\theta$, certificate & \textsf{CORRESPOND} $|$ \textsf{DIVERGE} $|$ \textsf{DECLINE}. \\
\typ{real}$^{+}$ & --- & A strictly positive real. \\
\typ{depth}    & --- & A non-negative partition depth (\Cref{sec:depth}). \\
\bottomrule
\end{tabular}
\end{center}

The subscript $\phi$ is a \emph{frame index}, assigned by the enclosing
\kw{under} block. Two types with different frame indices are different
types. This single device implements \Cref{rule:frames} and yields
\Cref{thm:scope}.

\subsection{Grammar}

We give the concrete grammar in EBNF. Nonterminals are
\textit{italic}; terminals are \texttt{typewriter}.

\begin{lstlisting}[style=grammar]
program     ::= { decl } { frame_block } report_stmt

decl        ::= "open" ident "=" string
              | "let" ident "=" expr
              | "method" ident "=" method_spec
              | "require" predicate

frame_block ::= "under" ident "{" { stmt } "}"

stmt        ::= binding | claim_stmt | sweep_stmt | for_stmt
              | relax_stmt | record_stmt | drop_stmt

binding     ::= "let" ident [":" type] "=" expr
              | "bind" ident "," ident "=" expr      # value, residue

expr        ::= project_expr | compare_expr | detect_expr
              | align_expr   | net_expr     | prim_expr

project_expr::= "project" expr "by" projector
projector   ::= "channels" "(" enc ")"
              | "spectral" "(" "coeffs" "=" int ")"
              | "contact"  "(" "medium" "=" num ")"
              | "cardinal" "(" ")"
enc         ::= "dna" | "protein"

compare_expr::= "compare" expr "against" expr "by" method_ref

detect_expr ::= "detect" detect_kind "in" expr "{" params "}"
detect_kind ::= "peaks" | "top"

align_expr  ::= "align" "central" "(" expr "," expr ")"
                        "response" "(" expr "," expr ")"
                        "under" corr_ref "{" params "}"

net_expr    ::= "perturb" expr "by" ident
              | "join" expr "and" expr
              | "unit" expr "anchors" int
              | "response" expr "by" ident

relax_stmt  ::= "relax" ident "until" "quiescent"
                "{" "eta" "=" num ";" "theta" "=" num "}"

sweep_stmt  ::= "sweep" ident "in" range "{" { stmt } "}"
range       ::= num ".." num "step" num | "[" num { "," num } "]"

for_stmt    ::= "for" ident "in" "items" "(" expr ")"
                [ "where" predicate ] "{" { stmt } "}"

claim_stmt  ::= "claim" string "=" expr
record_stmt ::= "record" ident { "," ident }
drop_stmt   ::= "drop" ident
report_stmt ::= "report" "to" string

params      ::= param { ";" param }
param       ::= ident "=" num
method_spec ::= "xcorr" "(" "normalised" ")"
              | "shader" "(" "cosine" ")"
              | "smith_waterman" "(" params ")"
              | "jaccard" "(" "k" "=" int ")"
\end{lstlisting}

\subsection{Frames and the \kw{under} block}
\label{sec:frames}

\begin{lstlisting}[style=nonum]
open wildtype = "wt.fa"
open variant  = "var.fa"

under nucleotide_frame {
    let a = project wildtype by channels(dna)
    let b = project variant  by channels(dna)
    let r = compare a against b by xcorr(normalised)
}

under protein_frame {
    let p = project wildtype by channels(protein)
    # let bad = compare a against p by xcorr(normalised)
    #   ^ scope error: `a` was created in nucleotide_frame
}
\end{lstlisting}

The commented line does not fail at runtime with a shape mismatch; it
fails at typecheck, because \typ{profile}$_{\texttt{nucleotide\_frame}}$
and \typ{profile}$_{\texttt{protein\_frame}}$ are distinct types. This is
\Cref{thm:scope}.

\subsection{Units and residue}
\label{sec:units}

Every comparison value is a pair. The user may bind both:

\begin{lstlisting}[style=nonum]
under N {
    bind score, res = compare q against t by xcorr(normalised)
    # `res` is in scope and MUST be consumed:
    record res
}
\end{lstlisting}

Leaving \texttt{res} unconsumed is a compile error. The user may
abandon it deliberately:

\begin{lstlisting}[style=nonum]
    drop res      # recorded in the report as an explicit abandonment
\end{lstlisting}

\noindent which typechecks, and prints in the report as a line reading
\texttt{residue abandoned at line \emph{n}}. There is no silent path.

\subsection{Relaxation}
\label{sec:relax}

\begin{lstlisting}[style=nonum]
under N {
    let d = compare unitA against unitB by demand
    relax d until quiescent { eta = 0.005 ; theta = 0.01 }
}
\end{lstlisting}

The compiler checks that \texttt{eta} is a positive literal or a
provably-positive expression. If it is not, the program is rejected with

\begin{lstlisting}[style=grammar]
error: `relax` requires a step floor bounded below (Assumption: effective update).
  eta = 0.0 is not positive; no termination bound can be emitted.
\end{lstlisting}

\noindent If it is, the compiler emits into the report the discharged
obligation
\[
  \text{steps} \;\le\; \left\lceil \dem_0/\eta \right\rceil
  \;=\; \left\lceil 0.4130/0.005 \right\rceil \;=\; 83 ,
\]
computed from the value of $\dem_0$ available at that point. The
termination of the loop is thus a theorem the compiler states, not a
hope the user holds.

\subsection{Partition depth}
\label{sec:depth}

Some analyses need an additive measure of how finely a space has been
subdivided --- how many nested distinctions have been drawn. We include
one primitive for it.

\begin{definition}[Partition depth]
\label{def:depth}
If a space is subdivided successively into $k_1,k_2,\dots,k_m$ parts,
its \emph{partition depth} at base $b$ is
$\mathrm{depth} = \sum_{j=1}^{m}\log_b k_j$.
\end{definition}

\begin{proposition}[Additivity]
\label{prop:depthadd}
If two subdivisions are independent, the depth of their product is the
sum of their depths.
\end{proposition}
\begin{proof}
The product subdivision has parts $k_j k'_l$ over the product index, and
$\log_b(k_jk'_l)=\log_b k_j+\log_b k'_l$; summing over the product
index and factoring gives the claim.
\end{proof}

Depth is a \typ{depth}-typed value and cannot be added to a
\typ{real}$^+$. Its use is to state how deep a hierarchical address is
before it is used as a prefilter (\Cref{sec:prog-homology}).

\subsection{The typing judgement}

We write $\Gamma;\phi \vdash e : \tau$ for ``in environment $\Gamma$ and
frame $\phi$, expression $e$ has type $\tau$''. Selected rules:

\begin{align*}
&\textsc{(Open)} &&
\frac{}{\Gamma;\phi \vdash \kw{open}\ x = s \;:\; \typ{seq}}
\\[6pt]
&\textsc{(Proj-Ch)} &&
\frac{\Gamma;\phi\vdash e:\typ{seq}}
     {\Gamma;\phi\vdash \kw{project}\ e\ \kw{by}\ \texttt{channels}(c)
       \;:\; \typ{profile}_\phi(\flo,\Res),\ \Res\ge\flo>0}
\\[6pt]
&\textsc{(Cmp-Lag)} &&
\frac{\Gamma;\phi\vdash e_1:\typ{profile}_\phi \quad
      \Gamma;\phi\vdash e_2:\typ{profile}_\phi}
     {\Gamma;\phi\vdash \kw{compare}\ e_1\ \kw{against}\ e_2\ \kw{by}\ \texttt{xcorr}
       \;:\;\typ{profile}_\phi}
\\[6pt]
&\textsc{(Cmp-Rank)} &&
\frac{\Gamma;\phi\vdash e_1:\typ{coord}_\phi \quad
      \Gamma;\phi\vdash e_2:\typ{coord}_\phi^{\,N}}
     {\Gamma;\phi\vdash \kw{compare}\ e_1\ \kw{against}\ e_2\ \kw{by}\ \texttt{shader}
       \;:\;\typ{ranked}_\phi}
\\[6pt]
&\textsc{(Align)} &&
\frac{\begin{array}{c}
       \Gamma;\phi\vdash c_A,c_B : \typ{unit}_\phi \qquad
       \Gamma;\phi\vdash r_A,r_B : \typ{response}_\phi \\
       \Gamma;\phi\vdash \varphi^{\mathrm{c}},\varphi^{\mathrm{r}} : \typ{corr}_{\phi,\phi}
       \qquad \theta>0 \text{ literal}
      \end{array}}
     {\Gamma;\phi\vdash \kw{align}\ \kw{central}(c_A,c_B)\ \kw{response}(r_A,r_B)\ \cdots
       \;:\;\typ{verdict}}
\\[6pt]
&\textsc{(Relax)} &&
\frac{\Gamma;\phi\vdash d:\typ{real}^{+} \qquad \eta>0 \text{ provable}}
     {\Gamma;\phi\vdash \kw{relax}\ d\ \kw{until\ quiescent}\ \cdots
       \;:\;\typ{verdict}
       \ \ \big[\text{emit } \lceil d/\eta\rceil\big]}
\end{align*}

Note what has \emph{no} rule: there is no rule producing $\Res=0$; there
is no rule deriving \typ{verdict} from fewer than four columns; there is
no rule relating $\typ{profile}_\phi$ to $\typ{profile}_\psi$ for
$\phi\ne\psi$; and there is no elimination rule for \typ{seq} other than
projection. The refusals of \Cref{sec:meta} are proved by inspecting the
absence of these rules.

% =====================================================================
\section{Operational semantics}
\label{sec:semantics}
% =====================================================================

We give a small-step semantics \cite{plotkin2004structural} over
configurations $\langle \Sigma, \mathcal{R}, s\rangle$ where $\Sigma$ is
a store, $\mathcal{R}$ is the accumulating report, and $s$ is the
statement sequence.

\subsection{Projection}

Channelisation is the single primitive shared by all families.

\begin{definition}[Channelisation]
\label{def:channelise}
Given a sequence $s$ of length $L$ over an alphabet $\mathcal{A}$ and an
encoding $c$, $\texttt{channels}(c)$ produces a matrix
$X\in\mathbb{R}^{C\times L}$:
\begin{itemize}[leftmargin=1.4em]
\item $c=\texttt{dna}$: $C=4$, $X_{\alpha,n}=\mathbb{1}[s_n=\alpha]$ for
      $\alpha\in\{\mathrm{A,C,G,T}\}$, then each row is centred by
      subtracting its mean.
\item $c=\texttt{protein}$: $C=3$, the rows being hydropathy
      \cite{kyte1982simple}, residue volume
      \cite{zamyatnin1972protein} and net charge at neutral pH, each
      normalised to $[0,1]$ and then centred.
\item $c=\texttt{cardinal}$: $C=2$, mapping each base to a unit vector
      in a fixed cardinal frame; the complementary strand maps to the
      negated trajectory, which is the formal expression of
      \Cref{sec:notinseq}.
\end{itemize}
Centring is not cosmetic: it removes the composition-only component so
that subsequent correlations respond to arrangement rather than base
frequency.
\end{definition}

\begin{align*}
&\textsc{(E-Proj)} &&
\langle \Sigma,\mathcal{R},\ \kw{let}\ x = \kw{project}\ s\ \kw{by}\ \texttt{channels}(c)\rangle
\;\longrightarrow\;
\langle \Sigma[x\mapsto X_{c}(s)],\ \mathcal{R}\cup\{\text{proj}\},\ \cdot\rangle
\end{align*}

\subsection{The three comparison methods}

\begin{definition}[Normalised cross-correlation]
\label{def:xcorr}
For query channels $Q\in\mathbb{R}^{C\times L_q}$ and target channels
$T\in\mathbb{R}^{C\times L_t}$ with $L_q\le L_t$, the score at lag $k$ is
\[
  r[k]\;=\;
  \frac{\sum_{\gamma=1}^{C}\sum_{n=0}^{L_q-1} Q_{\gamma,n}\,T_{\gamma,n+k}}
       {\|Q\|_F \cdot \|T_{[k,k+L_q)}\|_F},
  \qquad k=0,\dots,L_t-L_q,
\]
with the convention that a zero denominator yields $r[k]=0$. The
numerator is computed by FFT in $O((L_q+L_t)\log(L_q+L_t))$ time
\cite{cooley1965algorithm}; the denominator uses a rolling window norm.
Channels are summed \emph{coherently}, so phase agreement across
channels reinforces and disagreement cancels --- discarding phase and
correlating magnitudes alone is a strictly weaker procedure.
$r[k]\in[-1,1]$ by Cauchy--Schwarz.
\end{definition}

\begin{definition}[Spectral embedding]
\label{def:spectral}
For channels $X\in\mathbb{R}^{C\times L}$ and a coefficient count $m$,
let $\hat X_\gamma$ be the real FFT of row $\gamma$. The embedding is
\[
  E \;=\; \frac{1}{\max(L,1)}\,
  \Big(\big|\hat X_{\gamma,j}\big|\Big)_{\gamma\in[C],\ j\in\{1,\dots,m\}}
  \;\in\;\mathbb{R}^{Cm},
\]
zero-padded if $L$ is short. The DC term $j=0$ is \emph{skipped}: it
encodes only mean composition, which centring has already removed and
which would otherwise dominate the norm. Division by $L$ makes
embeddings of different-length sequences comparable. Embeddings are then
$\ell_2$-normalised, so cosine similarity is a dot product.
\end{definition}

\begin{definition}[Shader distance and hierarchical address]
\label{def:shader}
For an $\ell_2$-normalised query $q$ and a normalised database matrix
$B\in\mathbb{R}^{N\times Cm}$, the distance vector is $d = 1 - Bq \in
[0,2]^N$. A \emph{hierarchical address} of depth $D$ for a point in
$[0,1]^3$ is obtained by axis-cyclic bisection: at step $j$ bisect along
axis $j\bmod 3$ and emit the bit. Two points sharing an address prefix
of length $D$ lie in a common cell, giving a prefilter in the spirit of
\cite{morton1966computer,bentley1975multidimensional}; candidates
surviving the prefilter are reranked exactly by Smith--Waterman
\cite{smith1981identification}.
\end{definition}

\begin{definition}[Contact projection]
\label{def:contactproj}
$\texttt{contact}(\texttt{medium}=w)$ maps a solved network --- a set of
species with pairwise contact costs --- to a contact structure by taking
the costs as edge weights and attaching the medium to every species at
weight $w$. Setting $w$ to the minimum observed cost makes the medium
the cheapest distinction available and realises the floor without
inventing a scale; any other choice must be written explicitly, and the
report prints it.
\end{definition}

\subsection{Verdict evaluation}

\begin{align*}
&\textsc{(E-Align)} &&
\frac{\begin{array}{c}
 \dem^{\mathrm{c}} = \dem_{A\to B}+\dem_{B\to A}\quad
 \dem^{\mathrm{r}} = \dem_{R_A\to R_B}+\dem_{R_B\to R_A}\\
 \text{no relaxation declined}
\end{array}}
{\kw{align}\cdots \longrightarrow
 \begin{cases}
 \textsf{CORRESPOND} & \dem^{\mathrm{c}}\le\theta \wedge \dem^{\mathrm{r}}\le\theta\\
 \textsf{DIVERGE} & \text{otherwise}
 \end{cases}}
\\[8pt]
&\textsc{(E-Decline)} &&
\frac{\text{some relaxation reached an undefined update}}
     {\kw{align}\cdots \longrightarrow \textsf{DECLINE}}
\end{align*}

\subsection{Cost model}
\label{sec:costmodel}

The language commits to complexity bounds so that a user can predict
cost from the program text.

\begin{center}
\begin{tabular}{@{}lll@{}}
\toprule
Operation & Cost & Note \\
\midrule
\texttt{channels} & $O(CL)$ & linear scan \\
\texttt{xcorr} & $O((L_q+L_t)\log(L_q+L_t))$ & FFT \cite{cooley1965algorithm} \\
\texttt{spectral} & $O(CL\log L)$ & per sequence \\
\texttt{shader} & $O(NCm)$ & one matrix--vector product \\
prefix prefilter & $O(N)$ then $O(|\text{cand}|)$ & depth-$D$ address \\
\texttt{smith\_waterman} & $O(L_aL_b)$ per candidate & exact rerank \\
$\sep(u\mid v)$ & polynomial & max-flow \cite{ford1956maximal,stoer1997simple} \\
all-pairs $\sep$ & $|\Ver|-1$ flows & Gomory--Hu \cite{gomory1961multi} \\
\texttt{relax} & $\le\lceil \dem_0/\eta\rceil$ steps & \Cref{thm:dichotomy} \\
\bottomrule
\end{tabular}
\end{center}

\begin{remark}[Exact cuts in validation]
A validation suite should compute minimum cuts by exhaustive
enumeration over separating subsets, not by an approximate max-flow
routine, so that agreement between a computed value and a predicted one
cannot be an artefact of two approximations meeting. This is
exponential and is affordable only on small structures; synopsis
therefore admits an \texttt{exact} annotation on \texttt{contact} that
selects enumeration and refuses structures above a stated size.
\end{remark}

% =====================================================================
\section{Metatheory}
\label{sec:meta}
% =====================================================================

We now prove the properties that justify the design. Throughout, a
\emph{well-typed program} is one derivable in the judgement of
\Cref{sec:lang}.

\subsection{Totality}

\begin{theorem}[Strong normalisation]
\label{thm:total}
Every well-typed synopsis program terminates.
\end{theorem}

\begin{proof}
We exhibit a well-founded measure. Assign to each statement sequence the
lexicographic pair $(\mathcal{I}, \mathcal{S})$ where $\mathcal{S}$ is
the number of statements remaining and $\mathcal{I}$ is the sum, over
all enclosing iteration constructs, of the remaining iteration count.

Every construct that can repeat has a count fixed before it begins:
\begin{itemize}[leftmargin=1.4em]
\item \kw{sweep}~$x$~\kw{in}~$a..b$~\kw{step}~$s$ has
      $\lfloor (b-a)/s\rfloor+1$ iterations, determined by three values
      evaluated on entry; the grammar admits no expression for $s$ that
      is not positive.
\item \kw{sweep}~$x$~\kw{in}~$[\,\cdot\,]$ has the literal list length.
\item \kw{for}~$u$~\kw{in items}$(e)$ ranges over
      $I=\Ver\setminus\{\med\}$, finite by \Cref{def:fwg}; the optional
      \kw{where} guard only removes elements.
\item \kw{relax} has at most $\lceil \dem_0/\eta\rceil$ iterations by
      \Cref{thm:dichotomy}, and the typing rule \textsc{(Relax)}
      admits the construct only when $\eta>0$ is established, so the
      bound is a positive integer computed on entry.
\end{itemize}
There is no other repeating construct: the grammar has no \kw{while}, no
user-defined function and hence no recursion, and no fixpoint operator.
Each step of evaluation strictly decreases $(\mathcal{I},\mathcal{S})$
in the lexicographic order on $\mathbb{N}^2$, which is well-founded.
Hence no infinite reduction sequence exists.
\end{proof}

\begin{corollary}[Cost is predictable]
For each program the compiler can compute a closed-form upper bound on
the number of primitive operations, since every iteration count is
available statically or at loop entry.
\end{corollary}

\subsection{Floor preservation}

\begin{theorem}[Floor preservation]
\label{thm:floorpres}
Let $e$ be a well-typed expression of a comparison type evaluating to a
value carrying $(\flo,\Res)$. Then $\Res\ge\flo>0$.
\end{theorem}

\begin{proof}
Induction on the typing derivation. The base cases are the projection
rules, whose conclusions carry $\Res\ge\flo>0$ because
\Cref{def:contactproj} produces a structure with strictly positive
weights (\Cref{def:fwg}) and \Cref{thm:floor} then applies.

For the inductive step we check each rule that combines values.
\textsc{(Cmp-Lag)} and \textsc{(Cmp-Rank)} produce residues that are
sums or minima of subterm residues, each $\ge\flo>0$ by hypothesis, and
a finite sum or minimum of quantities $\ge\flo$ is $\ge\flo$.
\kw{perturb} is non-negative by \Cref{def:perturbation}, so by
\Cref{prop:skeleton} it does not decrease any separation cost and does
not lower $\flo$ unless it acts on the minimising edge, in which case it
raises it. \kw{join} forms a disjoint union sharing the medium; the
floor of the union is the minimum of the two floors, still strictly
positive, and the typing rule for \kw{join} requires both operands to be
in the same frame so that the medium is shared. No rule subtracts
residues. Hence the invariant is preserved.
\end{proof}

\begin{theorem}[No zero residue]
\label{thm:nozero}
There is no well-typed synopsis expression whose value carries
$\Res=0$.
\end{theorem}
\begin{proof}
Immediate from \Cref{thm:floorpres}: $\Res\ge\flo>0$.
\end{proof}

\subsection{The refusal theorems}

These state that certain scientifically invalid claims are not
expressible. Each is proved by exhibiting the absence of a typing rule,
which is a statement about the fixed grammar of \Cref{sec:lang} and so
is a genuine theorem about the language rather than an observation about
a particular implementation.

\begin{theorem}[Exactness is inexpressible]
\label{thm:noexact}
No well-typed synopsis program denotes the proposition ``$u$ and $v$ are
identical'', in the sense that no expression evaluates to a value
asserting zero separation cost between two items.
\end{theorem}

\begin{proof}
Such an expression would have to produce either (a) a comparison value
with $\Res=0$, excluded by \Cref{thm:nozero}; or (b) a boolean equality
on \typ{seq}, for which no rule exists --- the only elimination form for
\typ{seq} is \textsc{(Proj-Ch)} and its variants, all of which produce
comparison types carrying residue; or (c) a numeric expression
reconstructing identity from base-level access, which requires an
indexing operator on \typ{seq} that the grammar does not contain
(\Cref{rule:opaque}). All three routes are closed.
\end{proof}

\begin{theorem}[Arity of verdicts]
\label{thm:arity}
No well-typed synopsis program derives a value of type \typ{verdict}
from fewer than two central and two response columns.
\end{theorem}
\begin{proof}
\textsc{(Align)} is the only rule with \typ{verdict} in its conclusion
apart from \textsc{(Relax)}, which itself requires an \kw{align}
expression as its subject. \textsc{(Align)} has four column premises,
and a derivation lacking any premise is not a derivation. The grammar
production \textit{align\_expr} likewise requires both a
\texttt{central} and a \texttt{response} clause, so an attempt to omit
one fails at parse time with an arity error rather than producing a
weaker but well-formed value.
\end{proof}

\begin{theorem}[Scope safety]
\label{thm:scope}
No well-typed synopsis program compares two values created in distinct
\kw{under} blocks.
\end{theorem}
\begin{proof}
Each \kw{under} block introduces a fresh frame index $\phi$, and every
value constructed inside it receives type $\tau_\phi$. The comparison
rules \textsc{(Cmp-Lag)}, \textsc{(Cmp-Rank)} and \textsc{(Align)}
require their premises to carry the \emph{same} index, and there is no
subsumption, coercion or subtyping rule relating $\tau_\phi$ and
$\tau_\psi$ for $\phi\ne\psi$. Hence no derivation exists.
\end{proof}

\begin{corollary}[Encoding safety]
A program cannot compare a nucleotide-channelised value with a
protein-channelised value, nor a value registered in one receiver with
one registered in another. By \Cref{cor:notnoise} the latter would be a
category error rather than a tolerable approximation.
\end{corollary}

\begin{theorem}[Parameter completeness]
\label{thm:params}
For every well-typed program, the set of numeric parameters capable of
changing any verdict appears in the program text and is recoverable by
the compiler.
\end{theorem}
\begin{proof}
By \Cref{rule:nodefault} the grammar productions for
\textit{detect\_expr}, \textit{align\_expr}, \textit{relax\_stmt} and
\textit{method\_spec} each require a \textit{params} block, and
\textit{param} requires a literal on the right-hand side. No production
permits omission. Hence every such parameter is a token in the source,
and the parser collects them.
\end{proof}

\begin{corollary}[Residue accountability]
\label{cor:resacc}
Every residue produced by a well-typed program is either bound and
consumed, or explicitly dropped with a recorded source location. There
is no silent discard.
\end{corollary}
\begin{proof}
The \kw{bind} form introduces both components into scope; the
well-formedness condition on block exit requires every bound identifier
of residue type to be consumed by \kw{record} or \kw{drop}. \kw{drop}
emits a report entry. A block leaving a residue identifier unconsumed
does not typecheck.
\end{proof}

\subsection{What the metatheory does not give}

Honesty requires stating the boundary.

Totality (\Cref{thm:total}) says programs halt; it says nothing about
whether they halt \emph{soon}. A \kw{sweep} over $10^9$ points is
well-typed. The cost model of \Cref{sec:costmodel} is the tool for that,
not the type system.

Floor preservation (\Cref{thm:floorpres}) says residues stay positive;
it does not say they are \emph{small}, and a large residue is a
scientific problem the language reports but does not solve.

The refusal theorems say certain claims cannot be \emph{written}; they
do not say the claims a user does write are \emph{true}. Synopsis rules
out a family of category errors. It cannot rule out a badly chosen
threshold, a mis-specified receiver, or --- most importantly --- the
open problem of \Cref{sec:open}.

% =====================================================================
\section{Worked programs}
\label{sec:programs}
% =====================================================================

We give one complete program per experiment family. Each is a full
\texttt{.syp} file; none is a fragment.

\subsection{Matched-filter motif detection}
\label{sec:prog-motif}

\emph{Question.} Where in a target does a known motif occur, and with
what confidence against a background?

\begin{lstlisting}[caption={\texttt{motif\_scan.syp}}]
# Locate a motif in a target by coherent multichannel matched filtering.
# All thresholds are explicit; none has a default.

open motif  = "motif.fa"
open target = "chr7_region.fa"

under nucleotide {

    let q = project motif  by channels(dna)
    let t = project target by channels(dna)

    # Lag-indexed: the result type is profile, not ranked.
    bind r, res_r = compare q against t by xcorr(normalised)

    # Peak calling. Every parameter is stated.
    let hits = detect peaks in r {
        z            = 4.0 ;   # z-score against local background
        min_distance = 30  ;   # greedy maximum suppression radius
        min_score    = 0.35
    }

    claim "motif occurs at least once above z=4" = hits
    record hits, res_r
}

report to "motif_scan.report"
\end{lstlisting}

Points of note. The residue \texttt{res\_r} is bound and recorded; had
it been ignored the program would not compile
(\Cref{cor:resacc}). \texttt{r} has type \typ{profile}$_{\texttt{nucleotide}}$
and could not be passed to a top-$K$ selector, which consumes
\typ{ranked} (\Cref{rule:three}). The three peak parameters are in the
source, so the report can print them (\Cref{thm:params}).

A sweep over the decision threshold, to show the verdict's sensitivity
rather than assert one value:

\begin{lstlisting}[style=nonum]
    sweep z0 in 2.0..6.0 step 0.5 {
        let h = detect peaks in r { z = z0 ; min_distance = 30 ; min_score = 0.35 }
        record h
    }
\end{lstlisting}

\noindent This terminates in exactly $9$ iterations, fixed on entry
(\Cref{thm:total}).

\subsection{Spectral homology search}
\label{sec:prog-homology}

\emph{Question.} Which entries of a database are plausible homologues of
a query, when sequence-level search is too slow or too insensitive?

\begin{lstlisting}[caption={\texttt{homology.syp}}]
# Two-stage search: cheap spectral prefilter, exact rerank.

open query = "query.fa"
open db    = "swissprot_subset.fa"

under residue_space {

    let e  = project query by spectral(coeffs = 8)
    let DB = project db    by spectral(coeffs = 8)

    # Entry-indexed: the result type is ranked.
    bind s, res_s = compare e against DB by shader(cosine)

    let cands = detect top in s {
        k      = 200 ;         # prefilter width
        depth  = 12            # hierarchical address prefix length
    }

    # Exact rerank of survivors only. Costs O(La*Lb) per candidate.
    bind final, res_f = compare e against cands by smith_waterman(
        match = 2 ; mismatch = -1 ; gap = -2
    )

    claim "top 20 by exact rerank" = final
    record final, res_s, res_f
}

report to "homology.report"
\end{lstlisting}

The two-stage shape is forced by the cost model: the shader stage is one
matrix--vector product, $O(NCm)$, and the exact stage is quadratic per
candidate, so the address depth is the parameter trading sensitivity for
work. It is written in the source rather than defaulted. Note also that
\texttt{coeffs = 8} appears twice and must match --- comparing
embeddings of different dimension is a type error, since the dimension
is part of \typ{coord}$_\phi$.

\subsection{Four-column functional adjudication}
\label{sec:prog-fourcol}

\emph{Question.} Do a wild-type unit and a variant unit do the same
thing, given that content comparison cannot answer this
(\Cref{thm:falsefriend})?

\begin{lstlisting}[caption={\texttt{adjudicate.syp}}]
# The central question: same function, not same sequence.

open wt  = "wildtype.solved"
open var = "variant.solved"

# Named perturbation maps. Rule: responses may not be anonymous,
# because the verdict's independence of this choice is OPEN.
method pi_edge_raise = perturb_incident(magnitude = 0.10)

under cell_receiver {

    let Nwt  = project wt  by contact(medium = 0.05)
    let Nvar = project var by contact(medium = 0.05)

    # Central columns: the units with their anchor/residue split.
    let cA = unit Nwt  anchors 3
    let cB = unit Nvar anchors 3

    # Response columns: what each unit does to the receiver.
    # `by pi_edge_raise` is mandatory -- the arbitrary choice is named.
    let rA = response cA by pi_edge_raise
    let rB = response cB by pi_edge_raise

    let phi_c = corr from cA to cB by species_name
    let phi_r = corr from rA to rB by species_name

    # FOUR columns. Dropping `response(...)` is an arity error.
    let v = align central(cA, cB) response(rA, rB)
            under phi_c, phi_r { theta = 0.01 }

    relax v until quiescent { eta = 0.005 ; theta = 0.01 }

    claim "variant is functionally equivalent to wild type" = v
    record v
}

report to "adjudicate.report"
\end{lstlisting}

\noindent Three things this program cannot do, by construction:

\begin{enumerate}[label=(\roman*),leftmargin=2.2em]
\item It cannot conclude equivalence from \texttt{cA} and \texttt{cB}
      alone --- deleting the \texttt{response} clause is an arity error
      (\Cref{thm:arity}), so the mistake of \Cref{thm:falsefriend}
      cannot be made.
\item It cannot hide which response was chosen --- \texttt{pi\_edge\_raise}
      is named and printed (\Cref{rule:namedresponse}).
\item It cannot loop forever --- \texttt{eta} is positive, so the
      compiler emits the bound (\Cref{thm:dichotomy}).
\end{enumerate}

\subsection{An annotation transfer}

The stencil use of the framework: given an annotated reference and an
unannotated subject, transfer annotations by correspondence rather than
by sequence identity.

\begin{lstlisting}[caption={\texttt{transfer.syp}}]
open ref     = "reference.solved"
open subject = "subject.solved"

method pi_uniform = perturb_incident(magnitude = 0.10)

under tissue_receiver {

    let Nr = project ref     by contact(medium = 0.05)
    let Ns = project subject by contact(medium = 0.05)

    # For every item of the reference, ask whether the subject has a
    # unit that corresponds under BOTH column pairs.
    for u in items(Nr) where separation(u) > 0.10 {

        let cA = unit u  anchors 3
        let cB = nearest_unit Ns to cA
        let rA = response cA by pi_uniform
        let rB = response cB by pi_uniform

        let phi_c = corr from cA to cB by anchor_match
        let phi_r = corr from rA to rB by anchor_match

        let v = align central(cA, cB) response(rA, rB)
                under phi_c, phi_r { theta = 0.02 }

        claim "annotation of u transfers to its correspondent" = v
        record v
    }
}

report to "transfer.report"
\end{lstlisting}

The \kw{for} loop ranges over a finite item set with a guard, so it
terminates (\Cref{thm:total}). The guard $\sep(u)>0.10$ excludes items
too close to the medium to individuate reliably --- a filter the theory
supplies (\Cref{thm:floor}) rather than a heuristic. Each iteration
produces a \typ{verdict}, and a transfer that comes back
\textsf{DECLINE} is recorded as such rather than silently omitted.

% =====================================================================
\section{The report}
\label{sec:report}
% =====================================================================

The terminal form of every program is \kw{report to} a path. The
report is a structured document, not a log, and its schema is fixed.

\begin{center}
\begin{tabular}{@{}lp{9.6cm}@{}}
\toprule
Section & Content \\
\midrule
\textsc{Provenance} & Source file hash, every \kw{open} path and its
  hash, language version. \\
\textsc{Frames} & Each \kw{under} block, its receiver and encoding, and
  which values were created in it. \\
\textsc{Parameters} & Every numeric parameter in the source, with its
  line number. Complete by \Cref{thm:params}. \\
\textsc{Claims} & Each \kw{claim} string with its verdict, the four
  demands, and $\theta$. \\
\textsc{Residues} & Every residue: recorded value, or an explicit
  abandonment with line number (\Cref{cor:resacc}). \\
\textsc{Certificates} & For each \kw{relax}: $\dem_0$, $\eta$, the
  emitted bound $\lceil \dem_0/\eta\rceil$, and the realised step count. \\
\textsc{Declines} & Every \textsf{DECLINE}, with the step at which the
  update became undefined. \\
\textsc{Assumptions} & Each \kw{require} assertion, verbatim, marked
  \emph{asserted by user, not verified}. \\
\bottomrule
\end{tabular}
\end{center}

Two features deserve comment.

\textbf{Declines are reported, not suppressed.} A run in which half the
comparisons decline is a legitimate scientific outcome --- it says the
comparison as posed does not decide --- and a format that discards them
would misrepresent the run as more conclusive than it was.

\textbf{Assumptions are marked unverified.} \kw{require} lets a user
assert, for instance, that an expansion is floor-stabilising
(\Cref{def:stabilising}). The language cannot check this: it is a claim
about structures not yet observed. It therefore records the assertion
and labels it, so a reader of the report knows exactly which load the
conclusion is placing on the author's judgement rather than on the
computation.

% =====================================================================
\section{The open problem, and what the language does about it}
\label{sec:open}
% =====================================================================

\subsection{The negative result}

We tested \Cref{ass:respind} directly. Four solved steady-state
biochemical networks --- glycolysis, the citric-acid cycle, oxidative
phosphorylation and an EGFR/MAPK cascade, together supplying $35$
species and $35$ contact edges --- were converted to contact structures
by \Cref{def:contactproj}. For every species with at least two incident
non-medium edges, two admissible responses were constructed by raising
one or the other incident edge by the same magnitude, and the resulting
cross-demands compared. If \Cref{ass:respind} held, the two demands
would agree within $\theta$.

They do not.

\begin{center}
\begin{tabular}{@{}lr@{}}
\toprule
Quantity & Value \\
\midrule
Comparisons within $\theta=0.01$ & $28.6\%$ of $35$ \\
Mean discrepancy & $0.070$ \\
Maximum discrepancy & $0.307$ \\
Per-network median (glycolysis) & $0.063$ \\
Per-network median (citric-acid cycle) & $0.046$ \\
Per-network median (oxidative phosphorylation) & $0.077$ \\
Per-network median (EGFR/MAPK) & $0.100$ \\
\bottomrule
\end{tabular}
\end{center}

The maximum exceeds $\theta$ by more than an order of magnitude, and
\emph{every} network has a median above $\theta$. The failure is
systematic, not driven by one pathological pathway or one outlying
species. By \Cref{prop:respind}, the four-column verdict on these
networks depends on which admissible response the analyst constructed.

\subsection{What is missing}

We want to be precise about the nature of the gap, because it is easy to
mistake for a tuning problem.

What is missing is \textbf{a rule fixing the map from a unit to the
perturbation it induces}. The theory of \Cref{sec:contact}--\Cref{sec:blind}
is indifferent to which perturbation a unit is taken to cause; it only
requires that one be supplied. Physically, a unit does induce a
definite perturbation, so the map exists; we do not have a principled
construction of it from the data available. Better thresholds do not
help --- lowering $\theta$ makes agreement rarer, not commoner. More
data does not help --- by \Cref{cor:notnoise} the discrepancy is not
sampling error. What would help is a derivation of the response map from
the same contact structure the central columns are drawn from, so that
it is determined rather than chosen.

\subsection{What synopsis does about it}

A language cannot solve an open problem. It can refuse to let one hide.

\Cref{rule:namedresponse} requires every response column to name its
perturbation map. There is no anonymous response, no inferred response
and no default. Consequently:

\begin{itemize}[leftmargin=1.4em]
\item every program's source records which choice was made;
\item the report's \textsc{Parameters} section prints it
      (\Cref{thm:params});
\item two runs differing only in this choice are visibly different
      programs, so a disagreement between them is attributable rather
      than mysterious;
\item and the natural robustness check --- sweep the response map and
      see whether the verdict is stable --- is a three-line
      \kw{sweep} that the language makes cheap to write:
\end{itemize}

\begin{lstlisting}[style=nonum]
    sweep m in [0.05, 0.10, 0.20, 0.40] {
        let pi_m = perturb_incident(magnitude = m)
        let rA = response cA by pi_m
        let rB = response cB by pi_m
        let v  = align central(cA,cB) response(rA,rB)
                 under phi_c, phi_r { theta = 0.01 }
        record v
    }
\end{lstlisting}

\noindent A report showing the verdict flipping across this sweep is a
report telling the truth about the state of the theory. We prefer that
to a report showing a single confident verdict obtained by a default
nobody chose.

% =====================================================================
\section{Related work}
\label{sec:related}
% =====================================================================

\textbf{Sequence comparison.} The dynamic-programming tradition
\cite{needleman1970general,smith1981identification} and its heuristic
descendants \cite{altschul1990basic,altschul1997gapped,buchfink2015fast,steinegger2017mmseqs2}
optimise an alignment score --- a content quantity. Profile methods
\cite{eddy2011accelerated,durbin1998biological} and substitution
matrices \cite{henikoff1992amino} enrich the content model without
leaving it. Sketching approaches \cite{ondov2016mash,jaccard1912distribution}
trade exactness for speed, again within content.
\Cref{cor:orth} applies to all of them: they compute one coordinate of
a two-dimensional quantity. This is not a criticism of their execution;
it is a statement about what they are measuring.

\textbf{Genomic signal processing.} Representing sequence as a numeric
signal has a long history \cite{voss1992evolution,anastassiou2001genomic},
and the channelisation of \Cref{def:channelise} follows it. Our
contribution here is not the representation but its \emph{typing}: the
frame index makes an encoding choice part of the type, so that the
common error of comparing across encodings is caught statically.

\textbf{Function prediction and annotation.} Large-scale evaluations
show function prediction to be substantially unsolved
\cite{radivojac2013large}, and misannotation to propagate through
databases \cite{schnoes2009annotation}. Structure prediction
\cite{jumper2021highly} has transformed the structural half of the
problem without settling the functional one. Ontology and pathway
resources \cite{ashburner2000gene,kanehisa2000kegg} supply the
vocabulary in which function is stated but not a criterion for deciding
sameness. Our four-column criterion is such a criterion, blind to
annotation by construction.

\textbf{Networks and perturbation.} Treating cellular function as a
network property is standard \cite{barabasi2004network}, and
genome-scale perturbation screens now supply the response measurements
the response columns need
\cite{costanzo2016global,dixit2016perturb,replogle2022mapping}. The
framework here is, in that light, a proposal for what to \emph{compute}
from such screens: not a similarity between expression profiles but a
cross-demand between contact structures.

\textbf{Graph cuts.} Minimum cuts and their duality with flows
\cite{ford1956maximal,gomory1961multi,stoer1997simple,karger1996new,schrijver2003combinatorial}
are the computational substrate. Normalised cuts in vision
\cite{shi2000normalized,boykov2004experimental} share the intuition that
a boundary is the meaningful object; spectral connectivity
\cite{fiedler1973algebraic} offers a relaxation we do not use because we
need exactness for validation.

\textbf{Domain-specific languages.} The case for a DSL over a library
\cite{hudak1996building,mernik2005when,fowler2010domain} is usually made
on expressiveness. Ours is made on refusal, which places it nearer to
dimension-typed numerics \cite{kennedy1997relational}, linear types
\cite{wadler1990linear} and refinement types
\cite{freeman1991refinement}, where the type system exists to make
certain wrong programs unwritable. Totality follows
\cite{turner1995total}. Biological DSLs
\cite{pedersen2009towards,lakin2011design} and exchange formats
\cite{hucka2003systems} model systems; synopsis adjudicates claims about
them, which is a different job. Workflow systems
\cite{koster2012snakemake,difrancesco2015nextflow} give reproducible
orchestration but are agnostic about the scientific validity of the
steps they orchestrate; the report of \Cref{sec:report} is closer in
spirit to provenance models \cite{moreau2011provenance,wilkinson2016fair}.

\textbf{Reproducibility.} The no-defaults rule is a direct response to
the analytic-flexibility literature
\cite{simmons2011false,gelman2013garden,ioannidis2005why,baker2016reproducibility,sandve2013ten}.
Our contribution is to move the discipline from a checklist into a
grammar.

% =====================================================================
\section{Limitations}
\label{sec:limits}
% =====================================================================

We list what we do not claim.

\begin{enumerate}[label=\textbf{L\arabic*.},leftmargin=2.6em]
\item \textbf{Response independence is open.} \Cref{sec:open}. This is
      the principal gap and it is not small: without a rule fixing the
      response map, the four-column verdict is analyst-dependent. Every
      other limitation here is minor by comparison.

\item \textbf{Contact structures must be supplied.} Synopsis consumes
      solved networks; it does not solve them. Where the contact costs
      come from --- kinetics, perturbation screens, inference --- is
      outside the language, and a bad structure yields a confident
      wrong verdict.

\item \textbf{Floor-stabilisation is asserted, not checked.}
      \Cref{def:stabilising} is a hypothesis about structures not yet
      observed. \kw{require} records it and marks it unverified; nothing
      verifies it.

\item \textbf{The anchor count $r$ is a free parameter.} \Cref{def:anchor}
      fixes the tie-breaking but not $r$. The language forces $r$ into
      the source and the report, which makes it visible, not principled.

\item \textbf{Totality bounds steps, not time.} \Cref{thm:total} and the
      remark in \Cref{sec:meta}. A well-typed program may still be
      infeasible.

\item \textbf{No generative capability.} By the decision of
      \Cref{sec:whylanguage}, synopsis certifies but does not design.
      Users wanting sequence design must generate elsewhere and certify
      here.

\item \textbf{Empirical scope is narrow.} The negative result of
      \Cref{sec:open} rests on four networks and $35$ species. It is
      enough to refute \Cref{ass:respind} --- refutation needs only
      counterexamples --- but not enough to characterise how the
      assumption fails.

\item \textbf{Opacity has a cost we assert rather than measure.}
      \Cref{sec:opaque} argues from three experiment families that
      \Cref{rule:opaque} loses nothing. A wider survey could falsify
      this, and the honest position is that we have checked three
      families, not all of them.
\end{enumerate}

% =====================================================================
\section{Conclusion}
\label{sec:conc}
% =====================================================================

The question practitioners ask of two genomic units is whether they do
the same thing. The vocabulary available answers a different question
--- whether they are written the same way --- and \Cref{cor:orth} shows
the substitution is not an approximation but a projection that loses a
dimension. Two units may be indistinguishable in content and unbounded
in effect (\Cref{thm:falsefriend}); two may be maximally divergent in
content and identical in effect (\Cref{cor:converse}).

Occupying the missing dimension requires a theory in which identity is a
cut rather than a label (\Cref{thm:invariant-cut}, \Cref{cor:labels}),
distinction has a strictly positive price (\Cref{thm:floor}), value is
receiver-relative (\Cref{thm:receiver}), and correspondence can be
certified blind (\Cref{def:crossdemand}) with guaranteed termination
(\Cref{thm:dichotomy}). It requires a verdict with four columns and
three values (\Cref{def:fourq}).

It also requires, we argue, a language rather than a library --- because
the results above are largely \emph{prohibitions}, and a prohibition
that lives in documentation is not enforced. Synopsis makes exactness
inexpressible (\Cref{thm:noexact}), single-column verdicts an arity
error (\Cref{thm:arity}), cross-frame comparison a scope error
(\Cref{thm:scope}), silent residue discard impossible
(\Cref{cor:resacc}), and every verdict-relevant parameter mandatory
(\Cref{thm:params}); and it terminates (\Cref{thm:total}) with the
bound printed.

The name records the ambition. The synoptic gospels are read in parallel
columns not because they agree word for word but because they align
sufficiently for the columns to be laid alongside, with the residue
displayed rather than reconciled away. That is what genomic comparison
should be doing, and \Cref{thm:form} says exactly why it is coherent to
do so: two structures may place no demand on each other while remaining
fully distinguishable. Sufficient alignment, not exact alignment.

We end where we are weakest. \Cref{ass:respind} fails on real networks
--- $28.6\%$ agreement, maximum discrepancy $0.307$ against
$\theta=0.01$ --- and we have no rule fixing the response map. We have
built the language so that this cannot be concealed: no response may be
anonymous, and the sweep that exposes an unstable verdict is three
lines. A language that makes its own open problem visible in every
program that touches it is, we think, the correct thing to build while
the problem remains open.

\bibliographystyle{plain}
\bibliography{references}

\end{document}
